% ==============================
% Dummy Overleaf Project (HE) — Linear Algebra Book
% Exact environments: definition, example, exercise, remark,
% theorem, lemma, propositionosition, corollary, proof, solution, hint
% Numbering by chapter; solution/hint unnumbered.
% ==============================
\documentclass[12pt]{report} 

% ---- Math & friends ----
\usepackage{amsmath,amssymb,mathtools}
\usepackage{amsthm}  
\allowdisplaybreaks
\usepackage{float}
\usepackage{hyperref}
\usepackage{graphicx}
\usepackage{booktabs}
\usepackage{enumitem}
\usepackage{todonotes}
\usepackage{braket}
\usepackage{array}

% ---- Hebrew + fonts ----
\usepackage{fontspec}
\usepackage{polyglossia}
\setdefaultlanguage{hebrew}
\setotherlanguage{english}
% Overleaf has Frank Ruehl CLM in TeX Live
\newfontfamily\hebrewfont[Script=Hebrew]{Frank Ruehl CLM}
\newfontfamily\hebrewfontsf[Script=Hebrew]{Frank Ruehl CLM}
\newfontfamily\hebrewfonttt[Script=Hebrew]{Frank Ruehl CLM}

% For bullets
\renewcommand{\labelitemi}{$\bullet$}   % first-level bullet
\renewcommand{\labelitemii}{$\circ$}    % second-level bullet (optional)

% ---- Theorem-like environments (by chapter) ----
\theoremstyle{plain}
\newtheorem{theorem}{משפט}[chapter]
\newtheorem{lemma}[theorem]{למה}
\newtheorem{proposition}[theorem]{טענה}
\newtheorem{corollary}[theorem]{מסקנה}

\theoremstyle{definition}
\newtheorem{definition}[theorem]{הגדרה}
\newtheorem{example}{דוגמה}[chapter]
\newtheorem{exercise}{תרגיל}[chapter]
\newtheorem{remark}{הערה}[chapter]

% Unnumbered for website folding
\newtheorem*{solution}{פתרון}
\newtheorem*{hint}{רמז}

% ---- Macros (optional) ----
\newcommand{\R}{\mathbb{R}}
\newcommand{\N}{\mathbb{N}}
\newcommand{\Z}{\mathbb{Z}}
\newcommand{\Q}{\mathbb{Q}}
\newcommand{\C}{\mathbb{C}}
\newcommand{\F}{\mathbb{F}}
\DeclareMathOperator{\rank}{rank}
\newcommand{\M}{\mathbb{M}}
\newcommand{\0}{\mathbf{0}}
\DeclareMathOperator{\Span}{Span}
\DeclareMathOperator{\Row}{Row}
\DeclareMathOperator{\Col}{Col}
\DeclareMathOperator{\Null}{N}
\DeclareMathOperator{\Ker}{Ker}
\DeclareMathOperator{\Image}{Im}
\DeclareMathOperator{\Id}{Id}

\title{ספר אלגברה לינארית — דגם אוברליף}
\author{צוות הקורס}
\date{\today}

\begin{document}
\maketitle

\chapter{נספח: שימושים בהנדסה ובמדעים}\label{ch:appendix}

גישת ספר זה היא כגישת הקורס: אלגברה לינארית כתחום מתמטי מופשט עם מושגים כמו וקטורים, מטריצות, העתקות לינאריות ולכסון (וגם גיאומטריה לא מופשטת). אולם עבור מהנדסים ומדענים, מושגים אלה הם גם שפה מעשית לתיאור מערכות מורכבות. כאשר מתארים מבנה, מכונה, תהליך כימי, רשת תחבורה, מאגר נתונים או מערכת ביולוגית, מופיעה שוב ושוב אותה תבנית בסיסית: קודם אוספים משתנים רבים ומתארים את הקשרים ביניהם, וכך מקבלים מודל שאפשר לתאר ע"י מטריצות.

במקרים רבים, הבעיה ההנדסית מובילה לממ"ל. הנה כמה דוגמאות: הכוחות במבנה צריכים לקיים משוואות שיווי משקל; הטמפרטורות בגוף מחומם צריכות לקיים משוואות מעבר חום; הזרימות ברשת צריכות לקיים שימור חומר. בכל המקרים האלה, מטריצה מרכזת את המידע על הקשרים בין המשתנים, ופתרון הממ"ל מאפשר לחשב גדלים שאי אפשר למדוד או להבין ישירות.

המטריצה אינה רק דרך נוחה לסדר מספרים. לעיתים קרובות היא מייצגת את המבנה הפנימי של הבעיה: מי משפיע על מי, אילו משתנים מחוברים זה לזה, מהם האילוצים, ומהי מידת הרגישות של המערכת לשינויים. בהנדסת בניין, למשל, מטריצת קשיחות מתארת כיצד כוחות גורמים לתזוזות. בניתוח נתונים, מטריצת נתונים ומטריצת שונות משותפת מתארות קשרים בין מדידות. בהנדסה כימית ובביוטכנולוגיה, מטריצות מתארות מאזני חומר ורשתות תגובה. אותה שפה מתמטית מופיעה בכל פעם בהקשר אחר.

גם ערכים עצמיים, וקטורים עצמיים ולכסון מקבלים משמעות יישומית ברורה. במבנים ובמכונות הם מתארים תדרים עצמיים וצורות תנודה; במודלים דינמיים הם קשורים ליציבות או לחוסר יציבות; ובניתוח נתונים הם מאפשרים לזהות כיוונים מרכזיים שבהם הנתונים משתנים. במובן זה, לכסון אינו רק פעולה חישובית, אלא דרך לבחור מערכת צירים טבעית יותר לבעיה.

מטרת נספח זה היא להראות כיצד הרעיונות שנלמדו בספר מופיעים בתחומים שונים של הנדסה ומדעים (הקשורים למגמות הלימוד שלכם). בכל תחום נציג בעיה אופיינית (לפעמים שתיים), נסביר היכן מופיעות מטריצות וממ"ליות, ונראה כיצד מושגים כמו הפיכות, תלות לינארית ולכסון מקבלים פירוש מעשי. בנוסף לתחומים אלה, נביא דוגמה לשימוש הקשור לבינה מלאכותית וחיפוש מידע.

הדוגמאות הן יחסית קצרות וכמובן לא מחליפות קורסים בתחומים הרלוונטיים, אבל אולי הן ישכנעו אתכם שיש תמורה לידע שצברתם באלגברה לינארית. זו אחת השפות המרכזיות שבהן מהנדסים מתרגמים מציאות מורכבת למודל שניתן לנתח, לשפר ולעשות בעזרתו חישובים. נציין שמשוואות דיפרנציאליות (רגילות וחלקיות) הן עוד שפה מרכזית כנ"ל שמתכתבת עם אלגברה לינארית, אך לא נעסוק בהן כאן.

\section{שיטת הריבועים הפחותים}

בסטטיסטיקה ובניתוח נתונים רוצים לעיתים להתאים מודל פשוט למדידות. למשל, ייתכן שנרצה לתאר קשר בין מספר שעות לימוד לבין ציון במבחן, או אולי קשר בין שטח דירה לבין מחירה. במקרים כאלה הנתונים אינם מתאימים בדרך כלל באופן מושלם למודל. לכן השאלה אינה תמיד האם קיימת התאמה מדויקת, אלא מהי ההתאמה הטובה ביותר.

שיטת הריבועים הפחותים היא קודם כל יישום מרכזי של אלגברה לינארית, כלומר אפשר להבין אותה לחלוטין במונחים של אלגברה לינארית לפי שעושים קישור לתחומים אחרים. שיטה זו מופיעה כאשר לממ"ל
$Av=w$
אין פתרון מדויק, ולכן מחפשים וקטור
$v$
שעבורו
$Av$
קרוב ככל האפשר ל-
$w$.
מבחינה גיאומטרית, מחפשים את ההיטל של
$w$
על מרחב העמודות של
$A$.
לכן השיטה נשענת ישירות על מושגים בסיסיים באלגברה לינארית: מרחב עמודות והטלה אורתוגונלית. את המושג האחרון לא הגדרנו בקורס (הוא חורג ממנו), אבל יש קשר ישיר לבסיס אורתוגונלי ותהליך גרם-שמידט. נעבור על נושא זה בקצרה כהקדמה לשיטת הריבועים הפחותים.

\subsection*{הטלה אורתוגונלית}

בהינתן 
$w\in\R^m$
ותת-מרחב
$V\subseteq\R^m$,
ההיטל האורתוגונלי של
$w$
על
$V$
מוגדר להיות הוקטור
$\text{Proj}_V(w)\in V$
שמקיים
$w-\text{Proj}_V(w)\perp V$.
נוח להשתמש בסימון המקוצר
$\hat{w}=\text{Proj}_V(w)$.
למעשה, תנאי האורתוגונליות 
$w-\hat{w}\perp V$
שקולים לממ"ל שיש לה פתרון יחיד. ניתן לראות זאת ע"י בחירת בסיס אורתונורמלי
$B=\Set{v_1,...,v_d}$
ל-
$V$
וכתיבת
$\hat{w}=\sum_{i=1}^d\alpha_iv_i$.
מכאן:
$$
\begin{aligned}
w-\sum_{i=1}^d \alpha_iv_i\perp V &\iff \forall 1\leq j\leq d : \left(w-\sum_{i=1}^d \alpha_iv_i\right)\cdot v_j=0 \\
& \iff \forall 1\leq j\leq d : \alpha_j=w\cdot v_j
\end{aligned}
$$
זה מראה שבהינתן בסיס אורתונורמלי, ההיטל נתון ע"י 
$$
.\hat{w}=\sum_{i=1}^d(w\cdot v_i)v_i
$$
נזכיר שיש נוסחה דומה עבור הוקטור המקורי, בהנחה שמשלימים את 
$B$
לבסיס אורתונורמלי לכל
$\R^m$
ע"י הוספת וקטורים
$v_{d+1},...,v_m$.
כך מקבלים
$$
\begin{aligned}
.w=\sum_{i=1}^m(w\cdot v_i)v_i, \qquad w-\hat{w}=\sum_{i=d+1}^m(w\cdot v_i)v_i
\end{aligned}
$$

\begin{figure}[H]
    \centering
    \includegraphics[width=0.6\textwidth]{images/Projection.png}
    \caption{ההיטל $\hat{w}\in\Col(A)$ הוא הוקטור עבורו $w-\hat{w}\perp \Col(A)$}
    \label{fig:Projection}
\end{figure}

\subsection*{היטל כפתרון אופטימלי לבעיית מרחק}

המשמעות הגיאומטרית של ההיטל
$\hat{w}$
היא שזו הנקודה הקרובה ביותר ב-
$V$
ל-
$w$.
הניסוח המתמטי הוא
$$
,\|\hat{w}-w\|=\min_{v\in V}\|v-w\|
$$
ובאופן שקול
$$
.\|\hat{w}-w\|^2=\min_{v\in V}\|v-w\|^2
$$
נוח להסתכל על ריבוע המרחק כי הוא שווה לסכום של ריבועים (של הפרשי הקוארדינטות), ומכאן שם השיטה 
\textbf{ריבועים פחותים}.
מדוע ההיטל הוא אכן הנקודה הקרובה ביותר? לכל
$v\in V$
מתקיים
$v-\hat{w}\in V$
מתכונות של תת-מרחב, וכידוע
$w-\hat{w}\perp V$
ולכן גם
$\hat{w}-w\perp V$.
מכאן, לפי משפט פיתגורס, נובע כי
$$
.\|v-w\|^2=\|(v-\hat{w})+(\hat{w}-w)\|^2=\|v-\hat{w}\|^2+\|\hat{w}-w\|^2\geq \|\hat{w}-w\|^2
$$
זהו אי-שוויון חזק לכל
$v\neq \hat{w}$
כי במקרה זה מתקיים
$\|v-\hat{w}\|>0$.
זה בדיוק מראה שיש נקודת מינימום יחידה לבעיית המרחק, והיא
$\hat{w}$.

\subsection*{ריבועים פחותים כשיטה לבניית מודל}

\begin{example}\label{ex:least-squares}
נניח שיש לנו שלוש מדידות
$(1,2),(2,3),(3,5)$.
המטרה היא להתאים להן ישר מהצורה
$y=ax+b$.
הפרמטרים הלא ידועים הם
$a$
ו-
$b$.

שלוש הנקודות אינן שייכות לישר אחד. דרך אחת לראות זאת היא ע"י ניסיון לפתור את הממ"ל שמתקבלת ע"י הצבת הנקודות במשוואת הישר:

$$
\begin{cases}
\begin{alignedat}{3}
 a &{}+{}& b &= 2\\
2a &{}+{}& b &= 3\\
3a &{}+{}& b &= 5
\end{alignedat}
\end{cases}
$$

דירוג מטריצת המקדמים המורחבת יראה שיש שורת סתירה, כלומר אין פתרון לממ"ל. זה מראה שאכן אין ישר שעובר בדיוק דרך כל שלוש הנקודות. מה עושים? במקום לחפש פתרון מדויק, מחפשים ישר שמקרב את הנתונים בצורה הטובה ביותר.

נכתוב את המערכת בצורה מטריציונית עם סימן קירוב
$\approx$
במקום שוויון מדויק. נקבל
$Av\approx w$
כאשר
$$
.A=
\begin{pmatrix}
1 & 1\\
2 & 1\\
3 & 1
\end{pmatrix},
\qquad
v=
\begin{pmatrix}
a\\
b
\end{pmatrix},
\qquad
w=
\begin{pmatrix}
2\\
3\\
5
\end{pmatrix}
$$

הוקטור
$Av$
הוא וקטור התחזיות של המודל המתאים לנקודות המדידה 
$1,2,3$.
כלומר, אם בוחרים ערכים מסוימים של
$a,b$
אז
$$
Av=
\begin{pmatrix}
a+b\\
2a+b\\
3a+b
\end{pmatrix}
$$
הוא הוקטור של הערכים שהישר חוזה בנקודות המדידה.
הערכים אלה אינם בהכרח הערכים 
$2,3,5$
שהתקבלו בפועל במדידות ומופיעים כקוארדינטות של הוקטור
$w$,
אבל הם אמורים להיות קירובים. נחזור לדוגמה זו בהמשך.

\end{example}

באופן כללי, נניח שיש 
$m$
מדידות הקשורות למשתנים
$x_1,...,x_n,y$. 
נניח בנוסף כי
$m>n$
ונחפש מודל לינארי מהצורה
$$
.y=a_0+a_1x_1+...+a_nx_n
$$
הפעם נסמן
$$
v=\begin{pmatrix}a_0\\\vdots\\a_n\end{pmatrix}
$$
ונקבל ממ"ל
$Av=w$,
כאשר
$A\in\M_{m\times (n+1)}(\R)$
ו-
$Av$
הוא וקטור התחזיות.
כל הוקטורים האפשריים מהצורה
$Av$
נמצאים במרחב העמודות של
$A$.
כלומר, מרחב העמודות של
$A$
הוא אוסף כל התחזיות שהמודל הלינארי יכול לייצר.
אם הוקטור
$w$
נמצא במרחב העמודות של
$A$,
אז קיימת התאמה מדויקת. לעומת זאת, אם
$w$
אינו נמצא במרחב העמודות, אז אין פתרון מדויק ולכן מחפשים את הוקטור ב-
$V=\Col(A)$
הקרוב ביותר ל-
$w$.
זהו ההיטל של
$w$
על
$V$,
כלומר
$\hat{w}=\text{Proj}_V(w)$.
דרך אחת למצוא אותו היא ע"י בחירת בסיס אורתונורמלי ל-
$V$,
אך זו לא הגישה היחידה והיא פחות נוחה פה כי אחר כך צריך לפתור את הממ"ל
$Av=\hat{w}$.
קיימת גישה שונה: מחפשים
$Av\in\Col(A)$
כך ש-
$w-Av$
מאונך למרחב העמודות, כלומר הוא מאונך לכל עמודה של
$A$
ובאופן שקול לכל שורה של
$A^T$.
מכאן שמתקיים
$$
.A^T(w-Av)=\underline{0}
$$

ע"י העברת אגפים מקבלים ממ"ל שנקראת מערכת משוואות נורמליות, שנתונה ע"י
\begin{equation}\label{eq:normal}
.A^TAv=A^Tw
\end{equation}
ניתן להוכיח שלכל 
$w\in\R^m$
לממ"ל
$A^TAv=A^Tw$
אכן קיים פתרון 
$\hat{v}\in\R^{n+1}$.
אם
$\rank(A)=n+1$,
הפתרון הנ"ל יחיד ו-
$A\hat{v}$
הוא בדיוק ההיטל 
$\hat{w}$.

\begin{example}
נחזור לדוגמה
\ref{ex:least-squares}.
בשביל אגף שמאל של המערכת הנורמלית, נחשב
$$
.A^TA=
\begin{pmatrix}
1 & 2 & 3\\
1 & 1 & 1
\end{pmatrix}
\begin{pmatrix}
1 & 1\\
2 & 1\\
3 & 1
\end{pmatrix}
=
\begin{pmatrix}
14 & 6\\
6 & 3
\end{pmatrix}
$$
ובשביל אגף ימין נחשב
$$
.A^Tw=
\begin{pmatrix}
1 & 2 & 3\\
1 & 1 & 1
\end{pmatrix}
\begin{pmatrix}
2\\
3\\
5
\end{pmatrix}
=
\begin{pmatrix}
23\\
10
\end{pmatrix}
$$

ולכן המערכת הנורמלית
\ref{eq:normal}
המתאימה לדוגמה זו נתונה ע"י

$$
.\begin{pmatrix}
14 & 6\\
6 & 3
\end{pmatrix}
\begin{pmatrix}
a\\
b
\end{pmatrix}
=
\begin{pmatrix}
23\\
10
\end{pmatrix}
$$

פתרון המערכת הוא
$a=\frac{3}{2},b=\frac{1}{3}$,
ולכן הישר המתאים בשיטת הריבועים הפחותים הוא
$$
.y=\frac{3}{2}x+\frac{1}{3}
$$

\begin{figure}[H]
    \centering
    \includegraphics[width=0.6\textwidth]{images/LeastSquares.png}
    \caption{הישר עבורו סכום ריבועי השגיאות מינימלי}
    \label{fig:LeastSquare}
\end{figure}

התחזיות של המודל מופיעות בוקטור

$$
.Av=
\begin{pmatrix}
\frac{3}{2}+\frac{1}{3}\\
3+\frac{1}{3}\\
\frac{9}{2}+\frac{1}{3}
\end{pmatrix}
=
\begin{pmatrix}
\frac{11}{6}\\
\frac{10}{3}\\
\frac{29}{6}
\end{pmatrix}
$$

לאחר שמצאנו את הישר
$y=\frac{3}{2}x+\frac{1}{3}$
אפשר להשתמש בו גם לתחזיות בנקודות שלא נעשו בהן מדידות (זו בעצם המטרה).
לדוגמה, עבור
$x=4$
המודל נותן
$$
.y=\frac{3}{2}\cdot 4+\frac{1}{3}=\frac{19}{3}
$$

תחזית כזו מבוססת על ההנחה שהקשר הלינארי שנמצא מתוך המדידות ממשיך לתאר בקירוב גם נקודות נוספות. כאשר התחזית נעשית מחוץ לתחום המדידות, זוהי אקסטרפולציה ויש לפרש אותה בזהירות.

\end{example}

הדוגמה לעיל מראה את ההבדל בין פתרון מערכת משוואות לבין התאמת מודל לנתונים.
בממ"ל רגילה מהצורה
$Av=w$
מחפשים פתרון מדויק.
בשיטת הריבועים הפחותים יודעים שלרוב אין פתרון מדויק, ולכן מחפשים את
$Av$
הקרוב ביותר ל-
$w$.
שיטת הריבועים הפחותים נותנת מודל, ולא רק התאמה לנקודות קיימות. לאחר שמצאנו את פרמטרי המודל, אפשר להציב בו ערכים חדשים ולקבל תחזיות שעשויות להיות טובות/סבירות בתחום מסוים.

\section{בינה מלאכותית: חיפוש סמנטי}

שימוש זה של אלגברה לינארית חורג ממגמות הלימוד שלכם, אבל יש לו השפעה רבה על זמינות הידע בחיים המודרניים. בחיפוש רגיל לפי מילים, המערכת מחפשת התאמה מילולית בין השאילתה לבין המסמכים או התמונות שבמאגר. בחיפוש סמנטי המטרה אחרת: למצוא פריטים הקרובים במשמעותם לשאילתה, גם אם הם אינם משתמשים באותן מילים בדיוק.

הרעיון המרכזי הוא לייצג טקסטים, תמונות, מוצרים או מסמכים כוקטורים במרחב רב-מימדי.
וקטור כזה נקרא
\textbf{שיבוץ}
.(embedding)
אם שני פריטים דומים במשמעותם, השיבוצים שלהם אמורים להיות קרובים זה לזה במרחב.

\begin{example}
השאילתות

\begin{center}
"מסעדה איטלקית בירושלים"
\end{center}

ו-

\begin{center}
"איפה יש פיצה בעיר ציון?"
\end{center}

אינן זהות מבחינה מילולית, אך הן קרובות מבחינה סמנטית. מערכת חיפוש סמנטי טובה אמורה לייצג אותן ע"י וקטורים קרובים. לעומת זאת, השאילתה

\begin{center}
"מתכון לעוגת אגסים"
\end{center}

אמורה לקבל וקטור רחוק יותר. הרי מי שרוצה לאפות עוגה כנראה לא מחפש מסעדה, ולהיפך.

\end{example}

\subsection*{כיצד נוצר שיבוץ של טקסט?}

במודל שפה מודרני, המשפט אינו מיוצג בדרך כלל ע"י רשימה פשוטה של מילים.
תחילה מפרקים את המשפט ליחידות קטנות, כגון מילים או חלקי מילים. כל יחידה מיוצגת ע"י וקטור. לאחר מכן המודל מעבד את כל הוקטורים יחד, כך שהייצוג של כל מילה מושפע מהקשר שבו היא מופיעה.

למשל, המילה "גשר" יכולה להופיע בהקשר של מבנה פיזי, אך גם בהקשר מופשט כמו "גשר בין תרבויות".
לכן הייצוג הוקטורי של המילה אינו תלוי רק במילה עצמה, אלא גם במילים שסביבה.

מבחינת אלגברה לינארית, המודל מבצע רצף של פעולות על וקטורים ומטריצות.
בכל שכבה של המודל מופיעות מכפלות של מטריצות בוקטורים, צירופים לינאריים של וקטורים, ונרמולים.
הפרטים של המודל עשויים להיות מורכבים, אך הרעיון החשוב עבורנו הוא פשוט:
המודל לומד תהליך חישובי שממיר טקסט לוקטור.

נסמן פונקציה זו באופן סכמטי ע"י
$E_{\text{text}}$.
אם השאילתה היא טקסט
$T$,
אז השיבוץ שלה הוא וקטור
$q=E_{\text{text}}(T)$
במרחב
$\R^d$.
חשוב להדגיש שהפונקציה
$E_{\text{text}}$
אינה ה"ל, בין היתר כי הטקסטים עצמם אינם איברים של מ"ו.
למרות זאת, בתוך החישוב מופיעות שוב ושוב פעולות לינאריות: כפל מטריצה בוקטור, צירופים לינאריים של וקטורים ומכפלות סקלריות. פעולות אלה משולבות עם פעולות לא לינאריות, ובסוף התהליך מתקבל שיבוץ וקטורי של הטקסט. המספר
$d$
הוא מימד השיבוץ. במערכות מעשיות הוא עשוי להיות מאות או אלפים.
למשל, אפשר לחשוב על
$d=768$
או
$d=1024$
כמימדים אופייניים. המספר המדויק תלוי במודל שבו משתמשים.

חשוב להדגיש שאין הכרח שלכל קואורדינטה של
$q$
תהיה משמעות פשוטה כמו "מסעדה", "איטלקית" או "ירושלים". המשמעות לא מסתתרת במספר צנוע של קוארדינטות, אלא בגיאומטריה של הוקטור כולו: וקטורים קרובים מייצגים טקסטים קרובים במשמעותם.

\subsection*{כיצד נוצר שיבוץ של תמונה?}

בתמונה המצב דומה, אך האובייקט הגולמי שונה. תמונה מורכבת מפיקסלים, וכל פיקסל מכיל ערכי צבע.
ייצוג ישיר של תמונה לפי הפיקסלים שלה אינו מתאים לחיפוש לפי משמעות. שתי תמונות של גשר יכולות להיות שונות מאוד בפיקסלים בגלל תאורה, זווית צילום, רקע או מרחק מהמצלמה.

לכן גם כאן משתמשים במודל שממיר את התמונה לוקטור.
נסמן את הפונקציה הרלוונטית ע"י
$E_{\text{image}}$.
אם התמונה היא
$I$,
אז השיבוץ שלה הוא
$v=E_{\text{image}}(I)$
כאשר
$v\in\mathbb{R}^d$.

שלב יצירת השיבוץ הוא שלב של עיבוד תמונה. המודל מקבל את הפיקסלים של התמונה, מעבד אותם, ומפיק מהם וקטור המתאר את התוכן החזותי של התמונה. עם זאת, עיבוד תמונה הוא תחום רחב יותר. לא כל עיבוד תמונה נועד ליצור שיבוץ, אך בחיפוש סמנטי מטרת העיבוד היא בדיוק זו: להמיר את התמונה לוקטור שאפשר להשוותו לוקטורים של טקסטים או תמונות אחרות.

במערכות מסוימות מאמנים את מודל הטקסט ואת מודל התמונה כך ששניהם יחזירו וקטורים באותו מרחב.
במצב כזה משפט ותמונה יכולים להיות מושווים זה לזה.
לדוגמה, השאילתה

\begin{center}
"מסעדה איטלקית בירושלים"
\end{center}

מקבלת וקטור
$q$
וכל תמונה במאגר מקבלת וקטור
$v_i$.
אם תמונה מסוימת אכן מייצגת מסעדה איטלקית בירושלים, נרצה שהוקטור שלה יהיה קרוב ל-
$q$.
כמובן, יש מספר מסעדות כאלו.

\subsection*{מאגר שיבוצים כמטריצה}

לאחר שכל הפריטים במאגר הומרו לוקטורים, אפשר לאסוף אותם למטריצה אחת.
נניח שבמאגר יש
$m$
פריטים, וכל פריט מיוצג ע"י וקטור באורך
$d$.
נסמן את הוקטורים ע"י
$v_1,v_2,\ldots,v_m$.

נארגן אותם כשורות של מטריצה

$$
.M=
\begin{pmatrix}
- & v_1^T & -\\
- & v_2^T & -\\
  & \vdots &  \\
- & v_m^T & -
\end{pmatrix}\in\M_{m\times d}(\R)
$$

כל שורה במטריצה היא שיבוץ של פריט אחד במאגר.

\begin{example}

אפשר לייצג מסמך שלם ע"י וקטור אחד, אך יש לזכור שזהו ייצוג דחוס של המסמך.
ייצוג כזה מתאים במיוחד למסמכים קצרים או לחיפוש לפי הנושא הכללי של המסמך. אם יש מיליון מסמכים קצרים וכל מסמך מיוצג ע"י וקטור באורך
$768$,
אז
$$
m=1{,}000{,}000,\qquad d=768
$$
ולכן
$M$
היא מטריצה מסדר
$1{,}000{,}000\times 768$. 
זוהי מטריצה גדולה, אך ניתן לפרש כל שורה כנקודה במרחב השיבוצים שמייצגים מסמכים.

\end{example}

\subsection*{כפל המטריצה בוקטור השאילתה}

נניח שהמשתמש כתב שאילתה, והמודל המיר אותה לוקטור
$q\in\R^d$.
כדי למצוא את הפריטים הדומים ביותר לשאילתה, מחשבים את הדמיון בין
$q$
לבין כל אחד מהוקטורים
$v_i$.

מה זה אומר? מקובל לנרמל את השיבוצים כך שאורכם הוא
$1$,
ולכן המכפלה הסקלרית
$v_i\cdot q$
מודדת את מידת הקרבה בין הפריט ה-
$i$
לבין השאילתה. נזכיר את הקשר לזווית: ככל שהמכפלה הסקלרית קרובה ל-
$1$,
הזווית 
$\theta$
בין הוקטורים קרובה ל-
$0$
כי מתקיים
$\cos\theta=v_i\cdot q$.
יש גם קשר למרחק בין הוקטורים, שהרי
$$
.\|v_i-q\|^2=2(1-v_i\cdot q)=2(1-\cos\theta)
$$

במקום לחשב את המכפלות הסקלריות לחוד, אפשר לחשב את כולן במסגרת אחת באמצעות כפל מטריצה בוקטור

$$
s=Mq
$$

וכך נקבל
$s\in\R^m$.
הקוארדינטה ה-
$i$
של
$s$
נתונה ע"י
$$
.s_i=v_i\cdot q
$$

המשמעות היא ש-
$s$
הוא וקטור ציוני הדמיון בין השאילתה לבין כל הפריטים במאגר.
לאחר מכן בוחרים את הפריטים שעבורם
$s_i$
הוא הגדול ביותר (קרוב ל-
$1$) 
כתוצאות החיפוש.

\subsection*{מדוע החישוב אמור להיות מהיר?}

במבט ראשון, החישוב נראה גדול מאוד.
אם
$M$
היא מטריצה בגודל
$1{,}000{,}000\times 768$,
אז החישוב
$Mq$
דורש
$768{,}000{,}000$
פעולות כפל ומספר דומה של פעולות חיבור. זהו מספר גדול, אך יש כמה סיבות לכך שהחישוב יכול להתבצע במהירות.
ראשית, כל שורה במטריצה אינה תלויה בשורות האחרות ולכן אפשר לחשב את המכפלות הסקלריות במקביל.
מעבדים מודרניים, ובעיקר מעבדים גרפיים, טובים במיוחד בביצוע הרבה מכפלות וחיבורים במקביל.
שנית, את מטריצת השיבוצים
$M$
מחשבים מראש.
כאשר פריט חדש נכנס למאגר, מחשבים את השיבוץ שלו ושומרים אותו. בזמן החיפוש אין צורך לחשב מחדש את כל השיבוצים, אלא רק את שיבוץ השאילתה.
שלישית, במאגרים גדולים מאוד לא תמיד בודקים את כל השורות של
$M$
אלא משתמשים בשיטות חיפוש מקורבות. הרעיון הוא לארגן את הוקטורים מראש במבנה נתונים שמאפשר למצוא במהירות מועמדים קרובים. כך המערכת בודקת רק חלק קטן מתוך המאגר.

\subsection*{סיכום}

חיפוש סמנטי הוא דוגמה מודרנית לשימוש באלגברה לינארית.
התחכום של הבינה המלאכותית נמצא ביצירת השיבוצים:
המודל לומד להמיר טקסטים ותמונות לוקטורים כך שקרבה גיאומטרית תשקף קרבה במשמעות. לצורך מציאת קרבה גיאומטרית קיים הכלי של מכפלה סקלרית, שאותו ניתן לשכלל ע"י חישוב מקבילי ושיטות חיפוש.

\section{הנדסת מכונות}

בהנדסת מכונות עוסקים בניתוח של גופים, מנגנונים וחלקים נעים.
גם כאשר הבעיה הפיזיקלית מורכבת, חלק מרכזי מהמודל מתבסס על רעיונות פשוטים מתוך אלגברה לינארית:
איסוף נעלמים לוקטור, כתיבת אילוצים בעזרת מטריצה, מעבר בין מערכות צירים, וזיהוי כיוונים שבהם התיאור הפיזיקלי נעשה פשוט יותר.

\subsection*{מטריצות סיבוב }

בהנדסת מכונות עוסקים לעיתים קרובות בגופים ובחלקים שמיקומם וכיוונם משתנים:
זרוע רובוטית, גלגל שיניים, להב של טורבינה, מנגנון צירי או חלק המורכב בזווית מסוימת בתוך מערכת גדולה יותר.
במצבים כאלה אותו וקטור פיזיקלי יכול להיכתב ביותר ממערכת צירים אחת.
למשל, אפשר לתאר נקודה ביחס לצירים המקומיים של החלק, או ביחס לצירים הקבועים של המכונה.

המעבר בין מערכות צירים הוא אחד המקומות הטבעיים שבהם מופיעות מטריצות בהנדסת מכונות.
במישור, סיבוב בזווית
$
\theta
$
(הפעם נמדוד אותה במעלות)
מתואר ע"י מטריצת הסיבוב
$$
.R_\theta =
\begin{pmatrix}
\cos\theta & -\sin\theta \\
\sin\theta & \cos\theta
\end{pmatrix}
$$

מטריצה זו מייצגת ה"ל: היא לוקחת וקטור במישור ומסובבת אותו בזווית
$\theta$
סביב הראשית.

נניח שלזרוע מכנית יש מערכת צירים מקומית המחוברת אליה.
אם נקודה מסוימת על הזרוע מתוארת במערכת הצירים המקומית ע"י הוקטור
$
p_{\mathrm{local}}
=
\begin{pmatrix}
a \\
b
\end{pmatrix}
$
והזרוע מסובבת בזווית
$\theta$
ביחס למערכת הצירים הקבועה, אז מיקום הנקודה במערכת הצירים הקבועה (גלובלית) הוא
$$
.p_{\mathrm{global}} = R_\theta p_{\mathrm{local}}
$$
כלומר, מטריצת הסיבוב מאפשרת לתרגם תיאור מקומי של נקודה לתיאור גלובלי שלה.

\begin{example}
נניח שנקודה נמצאת במרחק שתי יחידות לאורך הזרוע, ולכן במערכת הצירים המקומית שלה מתקיים
$$
.p_{\mathrm{local}}
=
\begin{pmatrix}
2 \\
0
\end{pmatrix}
$$

אם הזרוע מסובבת בזווית
$
30^\circ
$
ביחס לצירי המכונה, אז
$$
.R_{30^\circ}
=
\begin{pmatrix}
\cos 30^\circ & -\sin 30^\circ \\
\sin 30^\circ & \cos 30^\circ
\end{pmatrix}
=
\begin{pmatrix}
\frac{\sqrt{3}}{2} & -\frac{1}{2} \\
\frac{1}{2} & \frac{\sqrt{3}}{2}
\end{pmatrix}
$$

ולכן
$$
.p_{\mathrm{global}}
=
R_{30^\circ}
\begin{pmatrix}
2 \\
0
\end{pmatrix}
=
\begin{pmatrix}
\sqrt{3} \\
1
\end{pmatrix}
$$

כלומר, נקודה שנמצאת שתי יחידות קדימה ביחס לזרוע נמצאת, ביחס למערכת הצירים הקבועה, בנקודה
$(\sqrt{3},1)$.
דוגמה זו ממחישה שמטריצה אינה רק דרך לסדר מספרים.
היא יכולה לייצג פעולה גיאומטרית ממשית: סיבוב של וקטורים במישור.
\end{example}

כפל מטריצות מקבל משמעות ברורה כאשר מרכיבים כמה סיבובים בזה אחר זה.
נניח שחלק מכני אחד מסובב בזווית
$\theta$,
וחלק נוסף המחובר אליו מסובב בזווית נוספת
$\alpha$
ביחס אליו.
הסיבוב הכולל מתאר את ההרכבה של שתי הה"ל ומתקבל ע"י כפל מטריצות:
$R_\theta R_\alpha$.
במקרה של סיבובים במישור מתקיים (כצפוי וכפי שניתן לבדוק בעזרת זהויות טריגונומטריות)
$$
.R_\theta R_\alpha = R_{\theta+\alpha}
$$

בהקשר מכני אפשר לחשוב על זרוע רובוטית בעלת שני מפרקים:
כל מפרק מוסיף סיבוב, והמיקום הכולל של קצה הזרוע מתקבל מהרכבת הסיבובים.

לדוגמה, אם מפרק אחד מסובב את הזרוע בזווית
$20^\circ$
ומפרק נוסף מוסיף סיבוב של
$40^\circ$,
אז הסיבוב הכולל הוא
$R_{20^\circ}R_{40^\circ}=R_{60^\circ}$.

זוהי דוגמה פשוטה למשמעות של כפל מטריצות:
כאשר מטריצות מייצגות פעולות, הכפל שלהן מייצג ביצוע של פעולה אחת אחרי פעולה אחרת.

\subsection*{מטריצת אינרציה וצירים ראשיים}

לאחר שמבינים שמטריצות משמשות למעבר בין מערכות צירים, אפשר לראות שימוש נוסף של אלגברה לינארית בהנדסת מכונות:
מציאת מערכת צירים שבה תיאור של גוף נעשה פשוט יותר.

לגוף קשיח יש התנגדות לסיבוב, ותכונה זו תלויה בציר שסביבו מנסים לסובב את הגוף.
במישור, אפשר לתאר תכונה זו בעזרת מטריצה סימטרית מהצורה

$$
.I =
\begin{pmatrix}
I_{xx} & I_{xy} \\
I_{xy} & I_{yy}
\end{pmatrix}
$$

מטריצה זו נקראת \textbf{מטריצת האינרציה}.
האיברים שעל האלכסון מתארים מומנטים של אינרציה ביחס לצירים שנבחרו, והאיברים שמחוץ לאלכסון מתארים צימוד בין הצירים.

לעיתים הצירים שבחרנו לתיאור הגוף אינם הצירים הטבעיים ביותר שלו.
לכן מחפשים כיוונים מיוחדים שבהם פעולת מטריצת האינרציה פשוטה במיוחד.
מבחינה אלגברית, מחפשים וקטור
$v\neq\underline{0}$
וסקלר
$\lambda\in\R$
כך ש-
$I v = \lambda v$.
כלומר,
$v$
הוא ו"ע של מטריצת האינרציה, ו-
$\lambda$
הוא הע"ע המתאים. מאחר שמטריצת האינטרציה
$I$
סימטרית, ניתן למצוא בסיס אורתונורמלי של ו"ע לפי משפט הלכסון האורתוגונלי.

הפירוש ההנדסי הוא שהו"ע בבסיס זה הם \textbf{הצירים הראשיים} של הגוף, והע"ע הם \textbf{מומנטי האינרציה הראשיים}.
במערכת הצירים הראשית (לפי הבסיס של הצירים הראשיים) מטריצת האינרציה מקבלת צורה אלכסונית

$$
.D =
\begin{pmatrix}
I_1 & 0 \\
0 & I_2
\end{pmatrix}
$$

האפסים מחוץ לאלכסון מבטאים את העובדה שבמערכת צירים זו אין צימוד בין הכיוונים.
במילים אחרות, הלכסון מוצא מערכת צירים טבעית שבה תיאור תכונת האינרציה של הגוף נעשה פשוט יותר.

דוגמה זו מדגימה את המשמעות הגיאומטרית של ו"ע וע"ע.
הם אינם רק פתרונות של משוואה אלגברית, אלא כיוונים שבהם המערכת הפיזיקלית מתוארת בצורה הפשוטה ביותר.

\subsection*{סיכום}

בהנדסת מכונות אלגברה לינארית מופיעה באופן טבעי כאשר מתארים מיקום, כיוון ותכונות של גופים.
מטריצות סיבוב מאפשרות לעבור בין מערכות צירים שונות, וכפל מטריצות מתאר הרכבה של סיבובים.
בנוסף, לכסון מאפשר למצוא צירים ראשיים של גוף, שבהם מטריצת האינרציה מקבלת צורה אלכסונית ופשוטה יותר.

לכן אלגברה לינארית משמשת בהנדסת מכונות לא רק ככלי חישוב, אלא גם כשפה גיאומטרית לתיאור תנועה, כיוון ומבנה פנימי של גופים מכניים.

\section{הנדסת בניין}
\subsection*{מטריצת קשיחות}

אחת ההופעות המרכזיות של אלגברה לינארית בהנדסת בניין היא בניתוח התגובה של מבנים לעומסים.
כאשר מתכננים מבנה, גשר, מסגרת או רצפה, אחת השאלות הבסיסיות היא כיצד המבנה יזוז כאשר יופעלו עליו כוחות. במודלים לינאריים, התשובה לשאלה זו מתקבלת בעזרת ממ"ל.

נניח שמחלקים את המבנה למספר נקודות חיבור, הנקראות צמתים.
לכל צומת יש דרגות חופש, למשל תזוזה אופקית, תזוזה אנכית או סיבוב.
את כל התזוזות והסיבובים אוספים לוקטור 
$u$,
ואת כל הכוחות והמומנטים החיצוניים אוספים לוקטור
$f$.
הקשר בין הכוחות לבין התזוזות מתואר ע"י הממ"ל
$$
,K u = f
$$
כאשר המטריצה
$K$
נקראת
\textbf{מטריצת הקשיחות}
של המבנה. היא מתארת כיצד המבנה מתנגד לתזוזות.
אם רכיב מסוים במבנה קשיח מאוד, הוא תורם למטריצה איברים גדולים יותר.
אם שני צמתים מחוברים זה לזה, הקשר ביניהם מופיע בדרך כלל גם באיברים שמחוץ לאלכסון.

פתרון הממ"ל מאפשר לחשב את התזוזות של המבנה מתוך הכוחות הפועלים עליו. אם
$K$
הפיכה, אז לכל וקטור כוחות
$f$
קיים פתרון יחיד
$u=K^{-1}f$.
כלומר, הכוחות קובעים באופן חד-משמעי את תזוזות המבנה.
לעומת זאת, אם
$K$
לא הפיכה (המילה הנרדפת היא
\textbf{סינגולרית}),
אז קיים וקטור (למעשה וקטורים)
$v \neq \underline{0}$
כך ש-
$K v=\underline{0}$.

כל וקטור כזה מתאר תבנית תזוזה שאינה יוצרת כוחות מחזירים במודל.
במילים הנדסיות, זוהי תנועה אפשרית של המבנה שאינה מעוררת התנגדות קשיחותית.
תנועה כזו נקראת לעיתים \textbf{מנגנון}. 
המשמעות הפיזיקלית היא שהמבנה אינו מאולץ מספיק. למשל, אם גוף קשיח מונח במישור ללא קיבועים מתאימים, הוא יכול לזוז כולו ימינה או שמאלה בלי שאף איבר בו יתארך או יתכווץ. במקרה כזה אין כוח פנימי שמחזיר אותו למקומו.

\begin{example}
אפשר לראות זאת בדוגמה פשוטה של שני צמתים המחוברים בקפיץ אחד בלבד, ללא קיר קבוע.
נסמן את תזוזות הצמתים ב-
$u_1,u_2$.
לקפיץ חשוב רק ההפרש בין התזוזות, ולכן מטריצת הקשיחות היא מהצורה
$$
.K =
k
\begin{pmatrix}
1 & -1 \\
-1 & 1
\end{pmatrix}
$$

מטריצה זו אינה הפיכה. בפרט, מתקיים
$$
.K
\begin{pmatrix}
1 \\
1
\end{pmatrix}
=
\begin{pmatrix}
0 \\
0
\end{pmatrix}
$$

הוקטור
$
\begin{pmatrix}
1 \\
1
\end{pmatrix}
$
מתאר תזוזה שבה שני הצמתים זזים יחד באותה מידה.
במצב כזה אורך הקפיץ אינו משתנה, ולכן לא נוצר כוח מחזיר.
זוהי תנועה של הגוף כולו, ולא עיוות של הקפיץ.
\end{example}

הדוגמה לעיל ממחישה את הקשר בין סינגולריות (אי-הפיכות) של מטריצת הקשיחות לבין תנועות שאינן מלוות בהתנגדות של המבנה.
אם
$N(K) \neq \Set{\underline{0}}$,
אז קיימות תזוזות לא אפסיות שאינן יוצרות כוחות פנימיים.
בהנדסת מבנים, מצב כזה עשוי להעיד על כך שהמודל חסר קיבועים, שיש בו מנגנון תנועה, או שחלקים מסוימים במבנה אינם מחוברים באופן שמונע תנועה חופשית. לכן, בדיקת הפיכות של מטריצת הקשיחות אינה רק בדיקה אלגברית. היא עונה על שאלה הנדסית חשובה:
האם המבנה (לפחות לפי המודל) אכן מסוגל להתנגד לכל תזוזה אפשרית של דרגות החופש שלו? גם אם התשובה שלילית, ייתכן שהוספת תנאי שפה (קיבועים) "תתקן" את המטריצה כך שהמטריצה החדשה תהיה הפיכה.

עוד ביטוי לסינגולריות הוא בעצם כך שלא לכל וקטור כוחות
$f$
יש פתרון לממ"ל.
כאשר
$K$
סינגולרית, הממ"ל
$Ku=f$
פתירה אם ורק אם
$f\in\Col(K)$.
אם
$f\notin\Col(K)$,
אז אין וקטור תזוזות שיכול לאזן את הכוחות האלה במסגרת המודל.
גם לכך יש פירוש הנדסי:
העומסים שהופעלו אינם מתיישבים עם האילוצים והחיבורים של המבנה כפי שתוארו במודל.

\subsection*{מאמצים ראשיים וכיוונים עצמיים}

דוגמה נוספת לשימוש באלגברה לינארית בהנדסת בניין מופיעה בניתוח מאמצים.
בנקודה בתוך קורה, עמוד, לוח או אלמנט קרקע, מצב המאמצים במישור ניתן לתיאור בעזרת מטריצה סימטרית:

$$
\sigma =
\begin{pmatrix}
\sigma_x & \tau_{xy} \\
\tau_{xy} & \sigma_y
\end{pmatrix}
$$
כאשר
$\sigma_x$
ו-
$\sigma_y$
הם מאמצים נורמליים בכיווני הצירים, ו-
$\tau_{xy}$
הוא מאמץ גזירה.

הצירים שבהם מודדים את המאמצים נבחרים בדרך כלל מטעמי נוחות, למשל אופקי ואנכי.
אבל זו בחירה טבעית לנו, לא בהכרח לחומר עצמו.
לכן עולה שאלה: האם קיימים כיוונים מיוחדים שבהם תיאור המאמצים פשוט יותר?

התשובה ניתנת באמצעות לכסון. מחפשים ע"ע
$\lambda_1,\lambda_2\in\R$
ובהתאמה ו"ע
$v_1,v_2\neq\underline{0}$.

הפירוש ההנדסי הוא חשוב:
הו"ע של מטריצת המאמצים הם הכיוונים הראשיים של המאמץ, והע"ע הם המאמצים הראשיים.
בכיוונים אלה מאמץ הגזירה מתאפס, ולכן מצב המאמצים מתואר בצורה הפשוטה ביותר.

במילים אחרות, לכסון מטריצת המאמצים פירושו מעבר ממערכת צירים שרירותית למערכת צירים טבעית של החומר.
במערכת צירים זו מתקבלת מטריצה אלכסונית
$$
D=
\begin{pmatrix}
\sigma_1 & 0 \\
0 & \sigma_2
\end{pmatrix}
$$
כאשר
$\sigma_1, \sigma_2$
הם המאמצים הראשיים. האאפסים מחוץ לאלכסון מבטאים את העובדה שבכיוונים הראשיים אין מאמצי גזירה. נדגיש שזהו למעשה לכסון אורתוגונלי שמובטח לנו כי 
$\sigma$
היא מטריצה סימטרית. אם כן, מערכת הצירים החדשה נראית כמו שהיינו רוצים: שני צירים מאונכים.

זוהי דוגמה ישירה לכך שלכסון, עם דגש על לכסון אורתוגונלי, אינו רק כלי חישובי.
הוא מאפשר לזהות את הכיוונים שבהם הבעיה ההנדסית מקבלת את הצורה הפשוטה והטבעית ביותר.

\section{הנדסת תעשייה וניהול}

\subsection*{משאבים, אילוצים ותכנון ייצור}

בהנדסת תעשייה וניהול עוסקים לעיתים קרובות בתכנון מערכות שבהן יש לקבל החלטות תחת אילוצים.
לדוגמה, מפעל צריך להחליט כמה יחידות לייצר מכל מוצר, מחסן צריך להחליט כיצד לחלק מלאי בין סניפים, ומערכת שירות צריכה להקצות עובדים למשימות שונות.
בכל המקרים האלה מופיעים משתנים רבים, קשרים בין המשתנים, ומגבלות/אילוצים על משאבים.

אלגברה לינארית מספקת שפה נוחה לתיאור מצבים כאלה.
המשתנים נאספים לוקטור, והקשרים ביניהם נכתבים בעזרת מטריצה.
כך מתקבל תיאור קצר וברור של מערכת שיש בה הרבה נתונים.

\begin{example}
נחשוב על מפעל המייצר שני מוצרים.
נסמן ב-
$x_1$
את מספר היחידות שייוצרו מן המוצר הראשון, וב-
$x_2$
את מספר היחידות שייוצרו מן המוצר השני. נסתכל על הוקטור
$x =
\begin{pmatrix}
x_1 \\
x_2
\end{pmatrix}$.
נניח שלייצור המוצרים דרושים שני משאבים: שעות עבודה וחומר גלם.
כל יחידה מהמוצר הראשון דורשת
$2$
שעות עבודה ו-
$3$
יחידות חומר גלם.
כל יחידה מהמוצר השני דורשת
$4$
שעות עבודה ו-
יחידת חומר גלם. את הנתונים האלה אפשר לארגן במטריצה

$$
.A =
\begin{pmatrix}
2 & 4 \\
3 & 1
\end{pmatrix}
$$

השורה הראשונה מתארת את צריכת שעות העבודה, והשורה השנייה מתארת את צריכת חומר הגלם.
העמודה הראשונה מתארת את המוצר הראשון, והעמודה השנייה מתארת את המוצר השני.

כפל המטריצה בוקטור ההחלטות נותן את סך המשאבים הדרושים, שהרי מתקיים

$$
.A x =
\begin{pmatrix}
2 & 4 \\
3 & 1
\end{pmatrix}
\begin{pmatrix}
x_1 \\
x_2
\end{pmatrix}
=
\begin{pmatrix}
2x_1+4x_2 \\
3x_1+x_2
\end{pmatrix}
$$

הרכיב הראשון הוא מספר שעות העבודה הדרושות, והרכיב השני הוא כמות חומר הגלם הדרושה. אם לרשות המפעל עומדות
$100$
שעות עבודה ו-
$90$
יחידות חומר גלם, אז הדרישה שלא לחרוג מהמשאבים פירושה

$$
A x \leq b
$$

כאשר
$b =
\begin{pmatrix}
100 \\
90
\end{pmatrix}$.
באופן מפורש, הדרישה היא
$$
.\begin{pmatrix}
2 & 4 \\
3 & 1
\end{pmatrix}
\begin{pmatrix}
x_1 \\
x_2
\end{pmatrix}
\leq
\begin{pmatrix}
100 \\
90
\end{pmatrix}
$$

אי-השוויון בין הוקטורים הוא סימון נוח (שלא השתמשנו בו בקורס כי לא עסקנו באי-שוויונות) שמשמעותו היא אי-שוויון בכל רכיב לחוד.
זו דרך לכתוב מערכת של שני אילוצים:
$$
\begin{cases}
\begin{alignedat}{3}
2x_1 &{}+{}& 4x_2 &{}& \leq{}& 100 \\
3x_1 &{}+{}&  x_2 &{}& \leq{}& 90
\end{alignedat}
\end{cases}
$$

העמודות של
$A$
מתארות את "טביעת המשאבים" של כל מוצר. העמודה
$\begin{pmatrix}
2 \\
3
\end{pmatrix}$
מתארת את המשאבים הדרושים ליחידה אחת מהמוצר הראשון, ואילו העמודה
$\begin{pmatrix}
4 \\
1\end{pmatrix}$
מתארת את המשאבים הדרושים ליחידה אחת מהמוצר השני. כאשר מייצרים
$x_1$
יחידות מהמוצר הראשון ו-
$x_2$
יחידות מהמוצר השני, סך המשאבים הדרושים הוא הצירוף הלינארי
$$
.x_1
\begin{pmatrix}
2 \\
3
\end{pmatrix}
+
x_2
\begin{pmatrix}
4 \\
1
\end{pmatrix}
$$

כלומר, תכנון הייצור הוא למעשה בחירה של צירוף לינארי של עמודות המטריצה.
הצירוף הזה צריך להישאר בתוך מגבלות המשאבים.

לעיתים רוצים גם לחשב בדיוק כמה לייצר אם מחליטים לנצל את כל המשאבים.
במקרה כזה מחליפים את האילוצים באי-שוויון בממ"ל
$A x = b$.
בדוגמה שלנו מתקבלת הממ"ל
$$
.\begin{pmatrix}
2 & 4 \\
3 & 1
\end{pmatrix}
\begin{pmatrix}
x_1 \\
x_2
\end{pmatrix}
=
\begin{pmatrix}
100 \\
90
\end{pmatrix}
$$

יש פתרון יחיד 
$x=\begin{pmatrix}26 \\ 12\end{pmatrix}$. 
זהו פתרון מעשי - שני המספרים טבעיים.

\end{example}

באופן כללי: מטריצה אחת יכולה לתאר מערכת שלמה של קשרים בין מוצרים, משאבים ואילוצים. במקום לכתוב כל קשר בנפרד, משתמשים בצורה המטריציונית
$A x \leq b$
כאשר הוקטור
$x$
מתאר את ההחלטה לגבי מספרי היחידות, המטריצה
$A$
מתארת את צריכת המשאבים, והוקטור
$b$
מתאר את הכמויות הזמינות. אם לממ"ל
$Ax=b$
יש פתרון, הוא מתאר כמויות ייצור שמנצלות בדיוק את כלל המשאבים הרלוונטיים לבעיה.
אם הפתרון אינו מתאים מבחינה מעשית, למשל אם אחד הרכיבים שלילי, פירוש הדבר הוא שהתכנון הזה אינו אפשרי ככמויות ייצור.

כאשר מחליפים את אילוצי המשאבים בממ"ל
$A x = b$
יש לזכור שלא כל פתרון אלגברי הוא בהכרח פתרון מעשי.
המערכת עשויה לתת וקטור
$x$
שמקיים את המשוואות בדיוק, אך אחד מרכיביו שלילי.
בהקשר של תכנון ייצור, רכיב שלילי פירושו "לייצר כמות שלילית" ממוצר מסוים, ולכן אין לו משמעות מעשית.
באופן דומה, ייתכן שהפתרון כולל מספרים שאינם שלמים, אף שבפועל אפשר לייצר רק יחידות שלמות.
לכן, בבעיות תכנון, פתרון הממ"ל
$Ax=b$
הוא רק שלב אחד בניתוח.
לאחר שמוצאים פתרון, יש לבדוק אם הוא עומד גם באילוצים המעשיים של הבעיה, כגון (במקרה של שני מוצרים)
$$
.x_1 \geq 0,
\qquad x_2 \geq 0
$$
אם הפתרון האלגברי אינו עומד באילוצים אלה, פירוש הדבר אינו שהאלגברה שגויה, אלא שהדרישה לנצל בדיוק את כל המשאבים אינה מתאימה למציאות של הבעיה.
מלכתחילה צריך לבחור
$b$
שיתאים לאפשרויות המעשיות של המערכת. פירוש הדבר שלא רק שצריך להתקיים
$b\in\Col(A)$
כדי שיהיה פתרון, אלא גם ש-
$b$
יתקבל כצירוף לינארי של עמודות
$A$
שבו הסקלרים הם מספרים טבעיים.

מה קורה אם הפתרון לממ"ל אינו יחיד? כאשר מספר המשוואות קטן ממספר המשתנים, יש פחות אילוצים ממספר רכיבי ההחלטה.
לכן, אם הממ"ל
$A x = b$
פתירה (במובן המעשי), בדרך כלל יישארו דרגות חופש בבחירת
$x$.
כלומר, יהיו כמה פתרונות אפשריים.
בהקשר ניהולי, פירוש הדבר הוא שהאילוצים לבדם אינם קובעים תוכנית אחת, ויש להוסיף קריטריון בחירה כגון מזעור עלות או הגדלת רווח.

\section{הנדסת מערכות מידע ונתונים}

בהנדסת מערכות מידע ונתונים עוסקים באיסוף, ארגון וניתוח של כמויות גדולות של מידע.
המידע יכול להגיע ממערכות עסקיות, אתרי אינטרנט, חיישנים, אפליקציות, מערכות המלצה, רשתות חברתיות או ניסויים.
אחד הרעיונות המרכזיים הוא שאפשר לתאר נתונים בעזרת וקטורים ומטריצות.

נניח שאוספים נתונים על מספר לקוחות.
עבור כל לקוח מודדים כמה תכונות, למשל גיל, מספר רכישות בחודש, סכום רכישה ממוצע וזמן שימוש באפליקציה. באופן כללי, כל לקוח יכול להיות מיוצג ע"י וקטור

$$
x =
\begin{pmatrix}
x_1 \\
x_2 \\
\vdots \\
x_n
\end{pmatrix}
$$

כאשר כל רכיב מתאר תכונה אחת של הלקוח. אם יש הרבה לקוחות, אפשר לאסוף את כל הוקטורים האלה למטריצה
$X$
(זהו הסימון המקובל בתחום).
אם יש
$m$
לקוחות, וכל לקוח מתואר ע"י
$n$
תכונות, אפשר לאסוף את כל הנתונים למטריצה

$$
X =
\begin{pmatrix}
x_{11} & x_{12} & \cdots & x_{1n} \\
x_{21} & x_{22} & \cdots & x_{2n} \\
\vdots & \vdots & \ddots & \vdots \\
x_{m1} & x_{m2} & \cdots & x_{mn}
\end{pmatrix}
$$

כאשר כל שורה במטריצה מייצגת לקוח אחד, וכל עמודה מייצגת תכונה אחת.
האיבר
$x_{ij}$
הוא הערך של התכונה ה-
$j$
עבור הלקוח ה-
$i$.

\begin{example}
נניח כי שלושה לקוחות מתוארים בעזרת שתי תכונות בלבד: מספר רכישות וסכום רכישה ממוצע. אז אפשר לכתוב

$$
.X =
\begin{pmatrix}
5 & 120 \\
2 & 80 \\
8 & 150
\end{pmatrix}
$$

השורה הראשונה מתארת לקוח שביצע
$5$
רכישות וסכום הרכישה הממוצע שלו הוא
$120$.
השורה השנייה מתארת את הלקוח הבא, וכן הלאה.
\end{example}

\begin{example}
מערכות המלצה הן דוגמה טבעית נוספת לשימוש במטריצות.
אפשר לבנות מטריצה שבה השורות מייצגות משתמשים, העמודות מייצגות מוצרים, והאיבר במקום
$(i,j)$
מתאר את הדירוג שנתן המשתמש
$i$
למוצר
$j$.

למשל:

$$
R =
\begin{pmatrix}
5 & 4 & 0 & 1 \\
4 & 0 & 0 & 1 \\
1 & 1 & 5 & 4 \\
0 & 1 & 4 & 5
\end{pmatrix}
$$

כאן האפסים יכולים לציין דירוגים חסרים, כלומר מוצרים שהמשתמש עדיין לא דירג.
המטרה של מערכת המלצה היא להשתמש בדירוגים הידועים כדי להעריך דירוגים חסרים, וכך להציע למשתמשים מוצרים שעשויים להתאים להם.
\end{example}

\subsection*{הורדת מימד}

במערכות מידע ונתונים מספר התכונות עשוי להיות גדול מאוד.
למשל, משתמש יכול להיות מתואר ע"י מאות מאפיינים, מסמך ע"י אלפי מילים, ותמונה ע"י אלפי פיקסלים.
במצב כזה קשה לראות את המבנה של הנתונים באופן ישיר.

אחת המטרות בניתוח נתונים היא למצוא תיאור פשוט יותר של הנתונים:
כלומר, להחליף מספר גדול של תכונות במספר קטן יותר של כיוונים מרכזיים.
מבחינת אלגברה לינארית, זו שאלה על מציאת כיוונים חשובים במרחב, שהם ו"ע המתאימים לע"ע גדולים יחסית.

רעיון זה מופיע, למשל, בשיטה הנקראת
\textbf{ניתוח רכיבים ראשיים}
או 
\textbf{PCA}
בראשי תיבות באנגלית.
בשיטה זו בונים מטריצה סימטרית המתארת את הקשרים בין התכונות.
מטריצה זו נקראת
\textbf{מטריצת שונות משותפת}.
הו"ע שלה מצביעים על כיוונים מרכזיים שבהם הנתונים משתנים, והע"ע מתארים עד כמה הנתונים משתנים בכל כיוון.

בהינתן מטריצת שונות משותפת
$C$,
מחפשים וקטורים
$v_i\neq\underline{0}$
וסקלרים
$\lambda_i$
כך ש-
$$
.Cv_i=\lambda_i v_i
$$
הו"ע
$v_1,...,v_n$
מתארים כיוונים מרכזיים של הנתונים. ע"ע גדול יותר מצביע על כיוון שבו השונות של הנתונים גדולה יותר.
לכן אפשר לבחור כמה כיוונים מרכזיים בלבד ולתאר בעזרתם את הנתונים בקירוב.

לכסון מטריצת השונות המשותפת אינו מוריד מימד בפני עצמו.
הלכסון רק מחליף את מערכת הצירים המקורית במערכת צירים חדשה, המותאמת למבנה הנתונים.
הו"ע של המטריצה מתארים כיוונים מרכזיים בנתונים, והע"ע מודדים את מידת השונות של הנתונים בכל כיוון.

הורדת מימד מתקבלת כאשר בוחרים רק חלק מהכיוונים האלה. אם
$\lambda_1 \geq \lambda_2 \geq \cdots \geq \lambda_n$
הם הע"ע המתאימים ל-
$v_1,v_2,\ldots,v_n$,
אז לכל
$k<n$
ניתן לבחור רק את
$v_1,\ldots,v_k$
ולתאר את הנתונים באמצעות ההיטל שלהם על תת-המרחב הנפרש ע"י ו"ע אלה. כלומר במקום ההצגה המלאה
$$
,v=\sum_{i=1}^n\frac{v\cdot v_i}{v_i \cdot v_i}v_i
$$
מסתפקים בהיטל (שלא הגדרנו בקורס)
$$
.\text{Proj}_k(v)=\sum_{i=1}^k\frac{v\cdot v_i}{v_i \cdot v_i}v_i\in\Span(v_1,...,v_k)
$$

מכאן נוח להשתמש בוקטור הקוארדינטות המתאים להיטל לעיל לפי הבסיס 
$B_k=\Set{v_1,...,v_k}$,
כלומר הוקטור
$$
.[\text{Proj}_k(v)]_{B_k}=\begin{pmatrix}\frac{v\cdot v_1}{v_1 \cdot v_1}\\\vdots\\\frac{v\cdot v_k}{v_k \cdot v_k}\end{pmatrix}\in\R^k
$$
כך וקטור נתונים שתואר במקור בעזרת
$n$
קוארדינטות מתואר בקירוב בעזרת
$k$
קוארדינטות בלבד. הורדת מימד היא קירוב טוב כאשר רוב השונות של הנתונים מרוכזת ב-
$k$
כיוונים, או ליתר דיוק כאשר היחס
$$
\frac{\lambda_1+\cdots+\lambda_k}
{\lambda_1+\cdots+\lambda_n}
$$
קרוב ל-
$1$.
אם כן, הלכסון מספק את מערכת הצירים המתאימה, והבחירה לשמור רק את הכיוונים המרכזיים פירושה הורדת המימד.

\section{הנדסה כימית}

בהנדסה כימית עוסקים בתהליכים שבהם חומרים זורמים, מתערבבים, נפרדים או מגיבים.
גם כאשר התהליך עצמו מורכב, אחד הכלים הבסיסיים ביותר לניתוחו הוא מאזן חומר:
מה שנכנס למערכת, מה שיוצא ממנה, ומה שנוצר או נצרך בתוכה.

במקרים רבים, מאזני חומר מובילים לממ"ל.
המשתנים הם בדרך כלל ספיקות של זרמים, ריכוזים או כמויות של רכיבים שונים.
המשוואות מבטאות שימור חומר או קשרים בין זרמים.
לכן אלגברה לינארית מספקת דרך טבעית לארגן את הנתונים ולחשב את הגדלים הלא ידועים.

במערכות כימיות גדולות יותר יש לעיתים מספר רב של יחידות תהליך:
מערבלים, מפרידים, ריאקטורים, מחליפי חום וזרמי מחזור.
במקום לכתוב כל משוואה בנפרד, אפשר לאסוף את כולן לממ"ל
$Ax=b$.
כאן
$x$
הוא וקטור הנעלמים, למשל ספיקות או ריכוזים לא ידועים.
המטריצה
$A$
מתארת את מאזני החומר ואת הקשרים בין הזרמים.
הוקטור
$b$
מתאר נתונים ידועים, כגון הזנות חיצוניות, דרישות ייצור או כמויות נתונות.

מנקודת מבט של אלגברה לינארית, מתעוררות שאלות טבעיות:
האם לממ"ל יש פתרון? האם הפתרון יחיד? האם הנתונים שניתנו מספיקים כדי לקבוע את כל הספיקות? האם יש משוואות מיותרות, כלומר יש תלות לינארית?

שאלות אלה מתורגמות בהנדסה כימית לשאלות על מידת ההגדרה של התהליך.
אם למערכת יש פתרון יחיד, אז הנתונים והמאזנים מספיקים כדי לקבוע את כל הנעלמים.
אם יש אינסוף פתרונות, ייתכן שחסרים נתונים, למשל ספיקה אחת שלא נמדדה או יחס ערבוב שלא נקבע.
אם אין פתרון, ייתכן שהנתונים אינם עקביים, למשל ריכוזים וספיקות שאינם יכולים להתקיים יחד עם שימור החומר.

לדוגמה, נניח שמערכת המאזנים היא
$Ax=b$
ומספר המשתנים גדול ממספר המשוואות הבלתי תלויות. במקרה כזה, אם המערכת פתירה, בדרך כלל יישארו דרגות חופש. 
בהקשר כימי, פירוש הדבר הוא שקיימות כמה דרכים שונות להפעיל את התהליך כך שיעמוד באותם מאזנים.
כדי לבחור פתרון אחד יש צורך בנתון נוסף או בשיקול תכנוני נוסף, כגון עלות, בטיחות, זמינות חומר גלם או איכות המוצר.

לעומת זאת, אם מוסיפים נתונים רבים מדי, ייתכן שהמערכת תהיה לא עקבית. לדוגמה, ייתכן שמדידה אחת של ספיקה ומדידה אחרת של ריכוז אינן מתאימות זו לזו. במקרה כזה מערכת המשוואות עשויה להיות חסרת פתרון.
מצב כזה אינו בהכרח כישלון של המודל המתמטי; לעיתים הוא דווקא מצביע על טעות מדידה, הנחה לא נכונה או איבוד חומר שלא נלקח בחשבון.

\subsection*{תגובות כימיות ומטריצת סטויכיומטריה}

דוגמה נוספת לשימוש במטריצות בהנדסה כימית מופיעה בתיאור תגובות כימיות.
במערכת שבה יש כמה חומרים וכמה תגובות, אפשר לארגן את מקדמי התגובות במטריצה אחת.
מטריצה זו נקראת \textbf{מטריצת סטויכיומטריה}.

\begin{example}
נניח שבמערכת מופיעים שלושה חומרים:
$A,\; B,\; C$
ותגובה אחת
$$
.A+B \longrightarrow C
$$

אפשר לכתוב את השינוי בכמויות החומרים בעזרת הוקטור
$$
.\begin{pmatrix}
-1 \\
-1 \\
1
\end{pmatrix}
$$
הסימנים השליליים מציינים חומרים שנצרכים בתגובה, והסימן החיובי מציין חומר שנוצר. אם יש כמה תגובות, כל תגובה מוסיפה עמודה למטריצה. אם, למשל, מוסיפים את התגובה
$$
,B \longrightarrow C
$$
מתקבלת המטריצה

$$
.S =
\begin{pmatrix}
-1 & 0 \\
-1 & -1 \\
1 & 1
\end{pmatrix}
$$

השורות מתאימות לחומרים, והעמודות מתאימות לתגובות.
האיבר במקום מסוים במטריצה מתאר כמה יחידות של חומר מסוים נוצרות או נצרכות בתגובה מסוימת.
\end{example}

באופן כללי, עבור 
$n$
תגובות תתקבל מטריצה
$S$
בעלת 
$n$
עמודות. נסתכל על וקטור
$$
,r =
\begin{pmatrix}
r_1 \\
\vdots \\
r_n
\end{pmatrix}
$$
שהוא וקטור המתאר את מידת ההתקדמות של כל תגובה. אז המכפלה
$Sr$
מתארת את השינוי הכולל בכמויות החומרים. כך מטריצה אחת מסכמת את הקשרים בין תגובות לבין חומרים.

גם כאן מופיעים רעיונות בסיסיים של אלגברה לינארית.
העמודות של
$S$
הן וקטורי שינוי אפשריים.
צירופים לינאריים של העמודות מתארים שינויים שאפשר לקבל ע"י שילוב של תגובות.
אם שינוי רצוי בכמויות החומרים אינו נמצא במרחב העמודות של
$S$,
אז אי אפשר לקבל אותו באמצעות התגובות שתוארו במודל.

מרחב הפתרונות של מטריצת סטויכיומטריה מקבל פירוש כימי ברור.
וקטור
$r\in\Null(S)$
מקיים
$Sr=\underline{0}$
ומייצג צירוף תגובות שאין לו שינוי נטו בהרכב החומרים. במילים אחרות, התגובות יכולות להתרחש, אך בסך הכול כל חומר שנצרך נוצר מחדש ע"י תגובות אחרות.

לסיכום: מרחב העמודות
$\Col(S)$
מתאר אילו שינויים אפשריים, והדרגה
$\rank(S)$
היא המימד של מרחב זה. לעומת זאת, מרחב הפתרונות של
$\Null(S)$
מתאר אילו צירופי תגובות אינם משנים דבר בסך הכול. המימד של מרחב זה הוא
$n-\rank(S)$
לפי משפט הדרגה.

\section{הנדסת חומרים}

\subsection*{כיוונים ראשיים של תכונות חומר}

בהנדסת חומרים עוסקים בקשר בין מבנה החומר לבין תכונותיו:
חוזק, קשיחות, הולכת חום, הולכה חשמלית, התפשטות תרמית ותכונות נוספות.
בחומרים מסוימים התכונה הנמדדת אינה זהה בכל הכיוונים.
לדוגמה, גביש או חומר מרוכב עשויים להיות קשיחים יותר בכיוון אחד ופחות קשיחים בכיוון אחר. במצבים כאלה אי אפשר לתאר את התכונה בעזרת מספר יחיד בלבד. במקום זאת משתמשים במטריצה.
המטריצה מתארת כיצד התכונה פועלת בכיוונים שונים, וכיצד רכיבים בכיוון אחד עשויים להיות קשורים לרכיבים בכיוון אחר.

ליתר דיוק, תכונה של חומר מתוארת ע"י מטריצה כאשר היא מקשרת בין גירוי וקטורי לבין תגובה וקטורית.
למשל, כיוון ועוצמה של שינוי טמפרטורה יכולים לגרום לשטף חום בכיוון מסוים.
אם הקשר הוא
$y=Ax$,
אז האיברים שעל האלכסון מתארים תגובה באותו כיוון, ואילו האיברים שמחוץ לאלכסון מתארים \textbf{צימוד}: גירוי בכיוון אחד יוצר תגובה גם בכיוון אחר.

נניח שתכונה מסוימת של חומר במישור מתוארת ע"י מטריצה סימטרית
$$
.A =
\begin{pmatrix}
a & b \\
b & c
\end{pmatrix}
$$

המטריצה יכולה לייצג תכונת הולכה, קשיחות אפקטיבית או תכונה כיוונית אחרת של החומר.
הסימטריה מבטאת הדדיות בין הכיוונים: הצימוד בין הכיוון הראשון לשני שווה לצימוד בין הכיוון השני לראשון.
העובדה שהמטריצה סימטרית חשובה מבחינה מתמטית:
למטריצה סימטרית יש ע"ע ממשיים, ואפשר לבחור לה ו"ע מאונכים. 

נניח שבוחרים כיוון במישור, המתואר ע"י וקטור
$v\neq\underline{0}$.
בדרך כלל, כאשר מפעילים את המטריצה על הוקטור, הכיוון של הוקטור משתנה.
כלומר,
$Av$
אינו בהכרח מקביל ל-
$v$.
אבל קיימים כיוונים מיוחדים (ו"ע) שבהם פעולת המטריצה אינה משנה את הכיוון.
בהקשר של הנדסת חומרים, הו"ע מתארים
\textbf{כיוונים ראשיים} של החומר.
אלה הכיוונים שבהם התכונה המתוארת ע"י המטריצה מקבלת צורה פשוטה במיוחד.
הע"ע מתארים את גודל התכונה בכיוונים הראשיים.

למשל, אם מטריצה מסוימת מתארת הולכה תרמית במישור, אז הע"ע הגדול יותר מתאים לכיוון שבו החומר מוליך חום בצורה חזקה יותר, והע"ע הקטן יותר מתאים לכיוון שבו ההולכה חלשה יותר.
אם המטריצה מתארת קשיחות, אז הע"ע מתארים קשיחויות ראשיות בכיוונים המתאימים.

מבחינת אלגברה לינארית, זהו בדיוק הרעיון של לכסון (למעשה אורתוגונלי).
למטריצה
$A$
יש בסיס של ו"ע, ולכן אפשר לעבור ממערכת הצירים המקורית למערכת צירים חדשה הבנויה מהו"ע.
במערכת הצירים החדשה המטריצה מקבלת צורה אלכסונית

$$
.D =
\begin{pmatrix}
\lambda_1 & 0 \\
0 & \lambda_2
\end{pmatrix}
$$

האפסים מחוץ לאלכסון מבטאים את העובדה שבכיוונים הראשיים אין צימוד בין הצירים.
כלומר, במערכת הצירים הראשית התכונה של החומר מתוארת ע"י שני מספרים בלבד:
$\lambda_1,\lambda_2$.

\begin{example}
נניח שתכונה מסוימת של חומר מתוארת ע"י המטריצה

$$
.A =
\begin{pmatrix}
5 & 2 \\
2 & 2
\end{pmatrix}
$$

כדי למצוא את הכיוונים הראשיים פותרים את המשוואה
$\det(A-\lambda I)=0$.
כך מקבלים
$$
,(5-\lambda)(2-\lambda)-4=0
$$
ומכאן
$$
.\lambda_1=1,
\qquad\lambda_2=6
$$

עבור
$\lambda_1=1$
אפשר לבחור ו"ע
$
v_1=
\begin{pmatrix}
-1 \\
2
\end{pmatrix}
$,
ועבור
$\lambda_2=6$
ניקח
$
v_2=
\begin{pmatrix}
2 \\
1
\end{pmatrix}
$,

שני הו"ע מאונכים זה לזה, כצפוי כי 
$A$
סימטרית והע"ע שונים. לכן הם יכולים לשמש כמערכת צירים חדשה.
במערכת הצירים הזאת המטריצה מקבלת את הצורה האלכסונית

$$
.D =
\begin{pmatrix}
6 & 0 \\
0 & 1
\end{pmatrix}
$$

המשמעות ההנדסית היא שהחומר אינו מתנהג באותו אופן בכל הכיוונים.
יש כיוון ראשי שבו התכונה חזקה יותר, וכיוון ראשי אחר שבו היא חלשה יותר.
הלכסון מאפשר לזהות את הכיוונים האלה ואת הערכים המתאימים להם.
\end{example}

בהנדסת חומרים, מטרת הלכסון היא לעבור ממערכת צירים שרירותית למערכת צירים ראשית של החומר.
במערכת המקורית, מטריצה המתארת תכונה כיוונית עשויה לכלול איברים מחוץ לאלכסון, המבטאים צימוד בין הכיוונים. לאחר לכסון, המטריצה מקבלת צורה אלכסונית, והצימוד נעלם.
הלכסון מאפשר לתאר את החומר בצורה פשוטה יותר, ולזהות באילו כיוונים התכונה חזקה או חלשה יותר.

מערכת הצירים הראשית אינה רק תוצאה חישובית.
היא מאפשרת לתאר את החומר לפי הכיוונים שבהם התכונה הנבדקת מתנהגת באופן הפשוט ביותר.
משתמשים במערכת הצירים הראשית כדי לפשט את החישוב, להפריד בין כיווני תגובה שונים, ולזהות את הכיוונים החשובים לתכנון הנדסי. במקום לתאר את החומר ביחס לצירים שנבחרו מטעמי נוחות, מתארים אותו ביחס לצירים המותאמים לתכונות הפנימיות שלו. אם יודעים שהחומר מוליך חום היטב בכיוון מסוים, או שהוא קשיח במיוחד בכיוון אחר, אפשר לנצל זאת בתכנון: לבחור את כיוון הסיבים בחומר מרוכב, למקם עומסים ביחס לכיווני החוזק, או לתכנן את כיוון ההנחה של שכבות החומר.

\section{ביוטכנולוגיה}
\subsection*{נתונים ביולוגיים כמטריצות}

בביוטכנולוגיה עוסקים לא רק בתהליכים כימיים ובתגובות ביולוגיות, אלא גם במדידה וניתוח של מערכות ביולוגיות מורכבות. לדוגמה, אפשר למדוד אילו גנים פעילים בתאים שונים, כיצד תאים מגיבים לתרופה, או כיצד משתנה הרכב חלבונים בתנאים שונים. בכל המקרים האלה מתקבלת כמות גדולה של נתונים, ואלגברה לינארית מספקת דרך טבעית לארגן ולנתח אותם.

נחשוב על ניסוי שבו מודדים את רמת הביטוי של כמה גנים במספר דגימות.
כל דגימה יכולה להיות תא, רקמה, חיידק, או מערכת ביולוגית שנבדקה בתנאי ניסוי מסוים.
אם מודדים
$n$
גנים, אז כל דגימה מתוארת ע"י וקטור
$$
x =
\begin{pmatrix}
x_1 \\
x_2 \\
\vdots \\
x_n
\end{pmatrix}\in\R^n
$$

כאשר
$x_j$
הוא רמת הביטוי של הגן ה-
$j$
בדגימה הנתונה. אם יש
$m$
דגימות, אפשר לאסוף את כל הנתונים למטריצה

$$
.D =
\begin{pmatrix}
d_{11} & d_{12} & \cdots & d_{1n} \\
d_{21} & d_{22} & \cdots & d_{2n} \\
\vdots & \vdots & \ddots & \vdots \\
d_{m1} & d_{m2} & \cdots & d_{mn}
\end{pmatrix}
$$

השורות של
$D$
מייצגות דגימות, והעמודות מייצגות גנים.
האיבר
$d_{ij}$
הוא רמת הביטוי של הגן ה-
$j$
בדגימה ה-
$i$.
בצורה זו, טבלת מדידות ביולוגית הופכת למטריצה.

היתרון של הייצוג המטריציוני הוא שאפשר להתייחס לכל דגימה כאל נקודה במרחב רב-מימדי. שתי דגימות קרובות זו לזו אם רמות הביטוי של הגנים בהן דומות. שתי דגימות רחוקות זו מזו אם דפוסי הביטוי שלהן שונים מאוד.
לכן שאלות ביולוגיות הופכות לשאלות גיאומטריות: אילו דגימות דומות זו לזו? אילו גנים משתנים יחד? האם טיפול מסוים משנה את הדגימות בכיוון ברור במרחב הנתונים?

\begin{example}
נניח ששתי דגימות מתוארות ע"י הוקטורים
$$
.u =
\begin{pmatrix}
2 \\
5 \\
1
\end{pmatrix},
\qquad
v =
\begin{pmatrix}
3 \\
4 \\
1
\end{pmatrix}
$$

הרכיבים מייצגים רמות ביטוי של שלושה גנים.
אפשר למדוד את המרחק בין הדגימות. כאן מתקיים

$$
.\|u-v\|=\sqrt{(-1)^2+1^2+0^2}=\sqrt{2}
$$

מרחק קטן מעיד על כך שהדגימות דומות יחסית בדפוס הביטוי שלהן, אך כאן הוא לא קטן במיוחד (אולי רק באופן יחסי ל-
$\|u\|,\|v\|$).

אפשרות אחרת היא למדוד דמיון באמצעות מכפלה סקלרית.
אם שני וקטורים מצביעים בכיוון דומה, הדבר עשוי להעיד על דפוס ביטוי דומה.
כדי להתמקד בכיוון ולא בגודל, משתמשים לעיתים בנוסחה

$$
\frac{u\cdot v}{\|u\|\|v\|}
$$

שמתארת את קוסינוס הזווית בין שני וקטורים הנתונים.
כאשר המדד קרוב ל-
$1$,
הרי שהזווית קרובה ל-
$0$
והדגימות דומות מבחינת דפוס הביטוי היחסי שלהן. זהו המקרה בדוגמה, שכן מתקיים

$$
.\frac{u\cdot v}{\|u\|\|v\|}=\frac{27}{\sqrt{30\cdot 26}}\approx 0.967
$$

\end{example}

מטריצות מאפשרות גם להשוות בין גנים. אם מתבוננים בעמודות של
$D$,
כל עמודה מתארת את רמות הביטוי של גן מסוים בכל הדגימות. שתי עמודות דומות מתארות שני גנים שמתנהגים באופן דומה לאורך הניסוי. למשל, ייתכן ששני גנים עולים ויורדים יחד בתגובה לטיפול מסוים.
במקרה כזה אפשר לומר שיש ביניהם דפוס משותף.

כך מתקבל פירוש חשוב לשורות ולעמודות של מטריצת הנתונים:
השורות מאפשרות להשוות בין דגימות, והעמודות מאפשרות להשוות בין גנים.
אותה מטריצה מתארת בו זמנית שתי נקודות מבט שונות על הניסוי.

לעיתים רוצים לתאר דגימה חדשה כצירוף של דפוסים מוכרים. לדוגמה, ייתכן שיש דפוס ביטוי אופייני לתאים בריאים ודפוס אחר לתאים שעברו טיפול. אם דגימה חדשה דומה בחלקה לדפוס הראשון ובחלקה לדפוס השני, אפשר לנסות לתאר אותה כצירוף לינארי (לפחות בקירוב) של שני וקטורים מייצגים. לשם כך, נסמן את שני הדפוסים ע"י
$p_1,p_2$,
ואת הדגימה החדשה ע"י
$y$.
אז אפשר לשאול האם קיימים סקלרים
$
\alpha,\beta\in\R
$
כך ש-
$y \approx \alpha p_1+\beta p_2$.

מבחינה ביולוגית, פירוש הדבר הוא שהדגימה החדשה מוסברת בקירוב ע"י שילוב של שני דפוסי ביטוי מוכרים.
מבחינת אלגברה לינארית, זו שאלה על צירופים לינאריים ועל קירוב וקטור באמצעות תת-מרחב הנפרש ע"י וקטורים נתונים.

באופן כללי, יש מספר כלשהו של דפוסים מוכרים והם לא בהכרח נותנים צירוף לינארי שמתאר בדיוק את הדגימה החדשה - אלא בקירוב. אם כן, כותבים ממ"ל מהצורה
$P c \approx y$
כאשר עמודות המטריצה
$P$
הן דפוסים מוכרים, והוקטור
$c$
מכיל את הסקלרים בצירוף הלינארי המתאים. הקירוב
$Pc \approx y$
פירושו, בדרך כלל, שבוחרים את
$c$
כך שהמרחק
$\|Pc-y\|$
יהיה קטן ככל האפשר.
מבחינה גיאומטרית, מחפשים את הוקטור בתת-המרחב הנפרש ע"י עמודות
$P$
הקרוב ביותר ל-
$y$
(כלומר ההיטל שלו).
הדרך המתמטית לעשות זאת נקראת
\textbf{שיטת הריבועים הפחותים}
ויש הסבר לגביה בסעיף הבא, אמנם בהקשר אחר.

\section{סטטיסטיקה וניתוח נתונים}\label{section:statistics}

\section*{איכות התאמה ומדד $R^2$}

שיטת הריבועים הפחותים מאפשרת למצוא מודל המתאים לנתונים בצורה מיטבית, במובן של מזעור סכום ריבועי השגיאות. אבל לאחר שמצאנו את המודל, עולה שאלה נוספת: עד כמה המודל באמת מסביר את הנתונים?

בסטטיסטיקה משתמשים לשם כך במדדים של איכות התאמה.
אחד המדדים השכיחים ביותר הוא
$R^2$
הנקרא מקדם ההסבר. נניח שיש לנו וקטור מדידות

$$
y=
\begin{pmatrix}
y_1\\
y_2\\
\vdots\\
y_m
\end{pmatrix}
$$

ומודל לינארי מהצורה
$Ac\approx y$.
לאחר פתרון בעיית הריבועים הפחותים מקבלים וקטור תחזיות
$\hat{y}=Ac$
וקטור השגיאה הוא
$e=y-\hat{y}$.

המטרה של הריבועים הפחותים היא לבחור את
$c$
כך שהאורך
$\|e\|$
יהיה קטן ככל האפשר.

\subsection*{השוואה למודל הממוצע}

כדי להבין אם המודל טוב, משווים אותו למודל פשוט מאוד:
מודל שמנבא תמיד את ממוצע המדידות.

נסמן את ממוצע הערכים של

$
y
$

על ידי

$
\overline{y}
$

כלומר

$
\overline{y}=\frac{1}{m}(y_1+y_2+\cdots+y_m)
$

וקטור התחזיות של מודל הממוצע הוא

$
\overline{y}\mathbf{1}
=
\begin{pmatrix}
\overline{y}\\
\overline{y}\\
\vdots\\
\overline{y}
\end{pmatrix}
$

כאשר

$
\mathbf{1}
=
\begin{pmatrix}
1\\
1\\
\vdots\\
1
\end{pmatrix}
$

הגודל

$
\|y-\overline{y}\mathbf{1}\|^2
$

מודד את סכום ריבועי הסטיות של הנתונים מן הממוצע.
זהו מדד לשונות הכוללת בנתונים.

לעומת זאת, הגודל

$
\|y-\hat{y}\|^2
$

מודד את סכום ריבועי השגיאות של המודל.
זהו החלק שהמודל לא הצליח להסביר.

לכן מגדירים

$
R^2
=
1-
\frac{\|y-\hat{y}\|^2}
{\|y-\overline{y}\mathbf{1}\|^2}
$

אם

$
R^2=0
$

אז המודל אינו משפר את התחזית לעומת שימוש בממוצע בלבד.
אם

$
R^2=1
$

אז המודל מתאים לנתונים באופן מושלם.
ערך כמו

$
R^2=0.8
$

פירושו שהמודל מסביר בערך

$
80\%
$

מן השונות בנתונים.

\subsection*{הפירוש הגיאומטרי}

כאשר המודל כולל איבר קבוע, כלומר כאשר אחת העמודות של

$
A
$

היא וקטור האחדות

$
\mathbf{1}
$

מתקבל פירוש גיאומטרי פשוט.

בריבועים פחותים, וקטור התחזיות

$
\hat{y}
$

הוא ההיטל של

$
y
$

על מרחב העמודות של

$
A
$

וקטור השגיאה

$
e=y-\hat{y}
$

מאונך למרחב העמודות של

$
A
$

ובפרט מאונך ל-

$
\mathbf{1}
$

כעת מפרקים את הסטייה הכוללת מן הממוצע:

$
y-\overline{y}\mathbf{1}
=
(\hat{y}-\overline{y}\mathbf{1})
+
(y-\hat{y})
$

הווקטור

$
\hat{y}-\overline{y}\mathbf{1}
$

הוא החלק שהמודל מסביר, והווקטור

$
y-\hat{y}
$

הוא השגיאה שנותרה.
שני הווקטורים האלה מאונכים זה לזה, ולכן לפי משפט פיתגורס:

$
\|y-\overline{y}\mathbf{1}\|^2
=
\|\hat{y}-\overline{y}\mathbf{1}\|^2
+
\|y-\hat{y}\|^2
$

כלומר:

\begin{center}
שונות כוללת = שונות מוסברת + שונות לא מוסברת.
\end{center}

מכאן אפשר לכתוב גם

$
R^2
=
\frac{\|\hat{y}-\overline{y}\mathbf{1}\|^2}
{\|y-\overline{y}\mathbf{1}\|^2}
$

ולכן

$
R^2
$

מודד את החלק היחסי של השונות הכוללת שהמודל הצליח להסביר.

\subsection*{דוגמה}

נחזור לדוגמה שבה התאימו ישר לשלוש נקודות:

$
(1,2),\qquad (2,3),\qquad (3,5)
$

מצאנו את הישר

$
y=\frac{3}{2}x+\frac{1}{3}
$

ולכן וקטור המדידות הוא

$
y=
\begin{pmatrix}
2\\
3\\
5
\end{pmatrix}
$

וקטור התחזיות הוא

$
\hat{y}=
\begin{pmatrix}
\frac{11}{6}\\
\frac{10}{3}\\
\frac{29}{6}
\end{pmatrix}
$

וקטור השגיאה הוא

$
e=y-\hat{y}
=
\begin{pmatrix}
\frac{1}{6}\\
-\frac{1}{3}\\
\frac{1}{6}
\end{pmatrix}
$

ולכן

$
\|e\|^2
=
\frac{1}{6}
$

ממוצע המדידות הוא

$
\overline{y}
=
\frac{2+3+5}{3}
=
\frac{10}{3}
$

ולכן

$
y-\overline{y}\mathbf{1}
=
\begin{pmatrix}
-\frac{4}{3}\\
-\frac{1}{3}\\
\frac{5}{3}
\end{pmatrix}
$

ומתקבל

$
\|y-\overline{y}\mathbf{1}\|^2
=
\frac{14}{3}
$

לכן

$
R^2
=
1-
\frac{\frac{1}{6}}{\frac{14}{3}}
=
1-\frac{1}{28}
=
\frac{27}{28}
$

כלומר,

$
R^2\approx 0.964
$

המשמעות היא שהישר שמצאנו מסביר כ-

$
96.4\%
$

מן השונות במדידות.

\subsection*{סיכום}

מדד

$
R^2
$

משווה בין שגיאת המודל לבין השגיאה שהייתה מתקבלת אילו היינו משתמשים רק בממוצע.
מבחינה אלגברית, המדד מבוסס על נורמות, הטלות ואורתוגונליות.
מבחינה סטטיסטית, הוא מודד איזה חלק מן השונות בנתונים מוסבר על ידי המודל.

כך מתקבל קשר טבעי בין ריבועים פחותים, גיאומטריה של הטלות וניתוח נתונים.



\section{מדעי כדור הארץ}

\subsection*{מדידות עקיפות ובעיות היפוך}

במדעי כדור הארץ עוסקים לעיתים במערכות שאי אפשר למדוד ישירות. למשל, קשה לדעת באופן ישיר מהו מבנה השכבות בעומק הקרקע, היכן קיימים אזורים צפופים יותר, או כיצד משתנה תכונה פיזיקלית מתחת לפני השטח.
במקום זאת מבצעים מדידות עקיפות: מודדים זמני הגעה של גלים, שינויים בשדה הכבידה, שדה מגנטי, טמפרטורה או ריכוזים בנקודות שונות.

הבעיה המרכזית היא לעיתים בעיית היפוך: מתוך המדידות מנסים לשחזר את המודל שיצר אותן. כלומר, במקום לחשב מה יהיו המדידות אם המודל ידוע, מנסים להסיק את המודל מתוך המדידות.

נניח שמחלקים אזור תת-קרקעי למספר תאים קטנים. בכל תא יש תכונה לא ידועה, למשל צפיפות, מהירות מעבר של גל, או ריכוז של חומר.
נסמן ב-
$x_j$
את הערך של התכונה בתא ה-
$j$,
ואת וקטור הערכים האלה ב-

$$
.x =
\begin{pmatrix}
x_1\\
x_2\\
\vdots\\
x_n
\end{pmatrix}
$$

כעת מבצעים מספר מדידות. כל מדידה אינה בודקת תא אחד בלבד, אלא מושפעת מכמה תאים. למשל, אם גל עובר במסלול מסוים דרך הקרקע, זמן ההגעה שלו מושפע מכל התאים שהמסלול עבר דרכם. אם מדידה מסוימת עוברת בעיקר דרך תא מסוים, התרומה של אותו תא למדידה תהיה גדולה יותר.

אפשר לתאר זאת בעזרת מטריצה
$A$
בה
השורה ה-
$i$
של המטריצה מתארת את המדידה ה-
$i$,
העמודה ה-
$j$
מתארת את התא ה-
$j$,
והאיבר
$a_{ij}$
מתאר עד כמה התא ה-
$j$
משפיע על המדידה ה-
$i$.
למשל, הוא יכול לייצג את אורך המסלול שעבר דרך התא, או משקל אחר המתאר את תרומת התא למדידה. אם וקטור המדידות הוא

$$
,b =
\begin{pmatrix}
b_1\\
b_2\\
\vdots\\
b_m
\end{pmatrix}
$$

אז המודל הלינארי מקבל את הצורה

$$
.Ax=b
$$

המטריצה
$A$
מתארת את גיאומטריית המדידה, הוקטור
$x$
מתאר את תכונות תת-הקרקע שרוצים לשחזר, והוקטור
$b$
מתאר את המדידות שנאספו.

\begin{example}
נניח שיש ארבעה תאים ושלוש מדידות.
ייתכן שמטריצת המדידה תהיה

$$
.A =
\begin{pmatrix}
1 & 1 & 0 & 0\\
0 & 1 & 1 & 0\\
0 & 0 & 1 & 1
\end{pmatrix}
$$

במודל זה, המדידה הראשונה מושפעת מהתאים הראשון והשני, המדידה השנייה מושפעת מהתאים השני והשלישי,
והמדידה השלישית מושפעת מהתאים השלישי והרביעי. אם

$$
,x =
\begin{pmatrix}
x_1\\
x_2\\
x_3\\
x_4
\end{pmatrix}
$$

אז

$$
.Ax =
\begin{pmatrix}
x_1+x_2\\
x_2+x_3\\
x_3+x_4
\end{pmatrix}
$$

כלומר, כל מדידה היא סכום של התרומות מהתאים שדרכם עברה. אם המדידות שנאספו הן

$$
,b =
\begin{pmatrix}
5\\
7\\
6
\end{pmatrix}
$$

אז מקבלים את הממ"ל

$$
\begin{pmatrix}
1 & 1 & 0 & 0\\
0 & 1 & 1 & 0\\
0 & 0 & 1 & 1
\end{pmatrix}
\begin{pmatrix}
x_1\\
x_2\\
x_3\\
x_4
\end{pmatrix}
=
\begin{pmatrix}
5\\
7\\
6
\end{pmatrix}
$$

בה יש פחות משוואות ממשתנים ולמעשה יש לה אינסוף פתרונות. מבחינה גיאופיזית, פירוש הדבר הוא שהמדידות אינן מספיקות כדי לקבוע באופן יחיד את כל הערכים בתאים.
ייתכנו כמה מבנים שונים של תת-הקרקע המייצרים אותן מדידות.

\end{example}

במקרים אחרים ייתכן שמספר המדידות גדול ממספר המשתנים, אך המדידות כוללות רעש או שגיאות מדידה.
אז בדרך כלל לא תהיה התאמה מדויקת, ולכן כותבים

$$
A x \approx b
$$

ומחפשים את
$x$
שעבורו
$Ax$
קרוב ככל האפשר ל-
$b$.
זהו שימוש טבעי בשיטת הריבועים הפחותים שכבר דנו בה בסעיף
\ref{section:statistics}.
הוקטור
$Ax$
הוא וקטור המדידות שהמודל חוזה, והוקטור
$b$
הוא וקטור המדידות בפועל. המטרה היא לבחור מודל תת-קרקעי כך שהתחזיות יהיו קרובות למדידות.

הדרגה של
$A$
מקבלת כאן משמעות חשובה.
אם הדרגה קטנה ממספר המשתנים, יש כיוונים במרחב המודלים שאינם ניתנים לזיהוי מתוך המדידות.
כלומר, אפשר לשנות את
$x$
בדרך מסוימת בלי לשנות את
$Ax$.
במקרה כזה קיימים מודלים שונים של תת-הקרקע שמייצרים אותן מדידות.
מבחינה אלגברית, אלה הם וקטורים ב-
$\Null(A)$.
אם
$v\in\Null(A)$,
אז
$A(x+v)=Ax$.
לכן המודלים
$x$
ו-
$x+v$
אינם ניתנים להבחנה באמצעות המדידות הנתונות.

זוהי אחת הסיבות לכך שבעיות היפוך במדעי כדור הארץ הן מאתגרות. לעיתים המדידות אינן מספיקות, או שהן רגישות רק לחלק מהתכונות של המודל. לכן יש צורך במדידות נוספות, בהנחות פיזיקליות, או במידע מוקדם על המערכת.

\section{קוגניציה}

\subsection*{מטריצות מעבר בתהליכי למידה}

חקר הקוגניציה עוסק, בין היתר, בלמידה, זיכרון, תפיסה וקבלת החלטות.
בסעיף זה לא ננסה לתאר את המוח עצמו בעזרת מטריצות.
במקום זאת נתמקד במודל פשוט של תהליך למידה:
מצב שבו אדם או קבוצת נבדקים עוברים בהדרגה בין כמה מצבי ידע או ביצוע.

נניח שבניסוי למידה נבדק יכול להימצא באחד משלושה מצבים: מצב
$S_1$
שבו הנבדק עדיין אינו יודע את הכלל;
מצב ביניים
$S_2$
שבו הנבדק מזהה את הכלל בחלק מהמקרים;
מצב
$S_3$
שבו הנבדק שולט בכלל.

בכל שלב של הניסוי הנבדק עשוי להישאר באותו מצב או לעבור למצב אחר.
את התפלגות הנבדקים בין המצבים נייצג בעזרת וקטורי עמודה
$$
x_k =
\begin{pmatrix}
a_{1,k}\\
a_{2,k}\\
a_{3,k}
\end{pmatrix}
$$
כאשר
$a_{i,k}$
הוא החלק של הנבדקים הנמצאים במצב
$S_i$
בשלב
$k$.
לכן הוקטור
$x_k$
אינו מייצגים אדם יחיד, אלא התפלגות של קבוצה בין המצבים האפשריים.

כדי לתאר את המעבר משלב לשלב משתמשים במטריצה שנקראת
\textbf{מטריצת מעבר}
(\underline{לא}
מהסוג שנלמד בקורס בין בסיסים). זו מטריצה
$A$
כך שהאיבר
$a_{ij}$
הוא ההסתברות לעבור מהמצב
$S_j$
אל המצב
$S_i$
בצעד אחד. אם כן, כל עמודה של
$A$
מתארת את כלל האפשרויות למעבר ממצב
$S_j$
למצב כזה או אחר,
ולכן סכום כל עמודה הוא
$1$.
המעבר משלב
$k$
לשלב
$k+1$
מתואר ע"י
$$
.x_{k+1}=A x_k
$$
במילים: כפל במטריצה
$A$
מתאר צעד אחד בתהליך הלמידה.

\begin{example}
נניח שמטריצת המעבר היא

$$
.A=
\begin{pmatrix}
0.7 & 0.1 & 0 \\
0.3 & 0.6 & 0.1 \\
0 & 0.3 & 0.9
\end{pmatrix}
$$

העמודה הראשונה מתארת מה קורה לנבדקים הנמצאים במצב
$S_1$
בשלב הנוכחי:
$70\%$
נשארים במצב
$S_1$
ו-
$30\%$
עוברים למצב
$S_2$.
העמודה השנייה מתארת מה קורה לנבדקים הנמצאים במצב
$S_2$
והעמודה השלישית מתארת מה קורה לנבדקים הנמצאים במצב
$S_3$.

נניח שבתחילת הניסוי כל הנבדקים נמצאים במצב הראשון. אז

$$
.x_0=
\begin{pmatrix}
1\\
0\\
0
\end{pmatrix}
$$

לאחר שלב אחד מתקבל

$$
.x_1=A x_0
=
\begin{pmatrix}
0.7 & 0.1 & 0 \\
0.3 & 0.6 & 0.1 \\
0 & 0.3 & 0.9
\end{pmatrix}
\begin{pmatrix}
1\\
0\\
0
\end{pmatrix}
=
\begin{pmatrix}
0.7\\
0.3\\
0
\end{pmatrix}
$$

כלומר, לאחר שלב אחד
$70\%$
מהנבדקים עדיין נמצאים במצב
$S_1$
ו-
$30\%$
עברו למצב
$S_2$.
לאחר שני שלבים מתקבל
$x_2=A x_1$,
ולכן
$$
.x_2=A(Ax_0)=A^2x_0
$$

\end{example}

באופן כללי, לאחר
$k$
שלבים מתקבל

$$
.x_k=A^k x_0
$$

מנקודת מבט קוגניטיבית, המודל מאפשר לשאול שאלות כמו:
כיצד משתנה התפלגות הנבדקים בין מצבי הידע לאורך זמן?
האם רוב הנבדקים עוברים בהדרגה למצב של שליטה בכלל?
האם מצב הביניים נשאר משמעותי גם לאחר כמה שלבים?
ומהי ההתפלגות הצפויה לאחר מספר נתון של צעדי למידה?

מנקודת מבט של אלגברה לינארית, וקטור המצב
$x_k$
הוא וקטור במרחב
$\R^3$
וכפל במטריצה
$A$
הוא ה"ל המעבירה את וקטור המצב לשלב הבא.
לכן תהליך הלמידה מתואר ע"י סדרת הוקטורים

$$
.x_0,Ax_0,A^2x_0,A^3x_0,...
$$

כאשר רוצים לתאר את התהליך לאחר מספר גדול של שלבים, יש לחשב את
$A^k$
ולכן מתקבלת בעיה שעסקנו בה בקורס: כיצד מחשבים חזקה גבוהה של מטריצה?
אם המטריצה
$A$
ניתנת ללכסון, כלומר קיימת מטריצה הפיכה
$P$
ומטריצה אלכסונית
$D$
כך ש-
$A=PDP^{-1}$.
אז החזקות של
$A$
מקבלות את הצורה הפשוטה
$$
.A^k = PD^kP^{-1}
$$

כזכור, קל מאוד לחשב חזקה של מטריצה אלכסונית. אם

$$
,D=
\begin{pmatrix}
\lambda_1 & 0 & 0\\
0 & \lambda_2 & 0\\
0 & 0 & \lambda_3
\end{pmatrix}
$$

אז

$$
.D^k=
\begin{pmatrix}
\lambda_1^k & 0 & 0\\
0 & \lambda_2^k & 0\\
0 & 0 & \lambda_3^k
\end{pmatrix}
$$

במודל המעבר בין מצבים, פירוש הדבר הוא שהו"ע מתארים כיוונים טבעיים של התהליך, והע"ע מתארים כיצד כל כיוון כזה משתנה מצעד לצעד. אם
$x_0$
הוא וקטור המצב ההתחלתי, אז
$$
.x_k=A^kx_0=PD^kP^{-1}x_0
$$

במטריצת מעבר תקינה, הע"ע חסומים. למעשה, כל ע"ע
$\lambda$
מקיים
$|\lambda|\leq 1$.
ניתן להוכיח כי 
$1$
הוא בהכרח ע"ע, והוא בדרך כלל מתאר רכיב יציב של התהליך.
רכיבים המתאימים לע"ע שעבורם
$|\lambda|<1$
דועכים כאשר מעלים את המטריצה בחזקות גבוהות, משום ש-
$\lambda^k \to 0$
כאשר
$k\to\infty$.
לכן, לכסון של מטריצת מעבר יכול לעזור להבין אילו רכיבים בתהליך נעלמים עם הזמן ואילו רכיבים נשארים יציבים.

\begin{example}
נחזור למטריצה
$$
.A=
\begin{pmatrix}
0.7 & 0.1 & 0 \\
0.3 & 0.6 & 0.1 \\
0 & 0.3 & 0.9
\end{pmatrix}
$$

הע"ע הם
$$
.\lambda_1=1,\qquad
\lambda_2=\frac{3}{5}+\frac{\sqrt{3}}{10},
\qquad
\lambda_3=\frac{3}{5}-\frac{\sqrt{3}}{10}
$$

שני הע"ע האחרונים קטנים מ-
$1$
בערך מוחלט, ולכן החזקות שלהם דועכות ככל ש-
$k$
גדל.
מכאן שלאחר מספר גדול של שלבים יישאר בעיקר הרכיב המתאים לע"ע
$1$.
ליתר דיוק, לכל רכיב של
$A^k$
יש גבול במובן של חדו"א, ומתקיים
$$
.A^k
\xrightarrow{k \to \infty}
\begin{pmatrix}
\frac{1}{13} & \frac{1}{13} & \frac{1}{13}\\
\frac{3}{13} & \frac{3}{13} & \frac{3}{13}\\
\frac{9}{13} & \frac{9}{13} & \frac{9}{13}
\end{pmatrix}
$$

כלומר, ללא תלות בהתפלגות ההתחלתית 
$x_0$,
מתקיים

$$
.A^kx_0\xrightarrow{k \to \infty}
\begin{pmatrix}
\frac{1}{13}\\
\frac{3}{13}\\
\frac{9}{13}
\end{pmatrix}
$$

המשמעות היא שהמודל מתכנס להתפלגות היציבה שמתוארת ע"י הוקטור לעיל, שהוא ו"ע המתאים לע"ע 
$1$
עם נרמול כך שסכום הקוארדינטות הוא 
$1$.

\end{example}

\end{document}